% Options for packages loaded elsewhere
\PassOptionsToPackage{unicode}{hyperref}
\PassOptionsToPackage{hyphens}{url}
%
\documentclass[
  10pt,
]{article}
\usepackage{amsmath,amssymb}
\usepackage{iftex}
\ifPDFTeX
  \usepackage[T1]{fontenc}
  \usepackage[utf8]{inputenc}
  \usepackage{textcomp} % provide euro and other symbols
\else % if luatex or xetex
  \usepackage{unicode-math} % this also loads fontspec
  \defaultfontfeatures{Scale=MatchLowercase}
  \defaultfontfeatures[\rmfamily]{Ligatures=TeX,Scale=1}
\fi
\usepackage{lmodern}
\ifPDFTeX\else
  % xetex/luatex font selection
\fi
% Use upquote if available, for straight quotes in verbatim environments
\IfFileExists{upquote.sty}{\usepackage{upquote}}{}
\IfFileExists{microtype.sty}{% use microtype if available
  \usepackage[]{microtype}
  \UseMicrotypeSet[protrusion]{basicmath} % disable protrusion for tt fonts
}{}
\makeatletter
\@ifundefined{KOMAClassName}{% if non-KOMA class
  \IfFileExists{parskip.sty}{%
    \usepackage{parskip}
  }{% else
    \setlength{\parindent}{0pt}
    \setlength{\parskip}{6pt plus 2pt minus 1pt}}
}{% if KOMA class
  \KOMAoptions{parskip=half}}
\makeatother
\usepackage{xcolor}
\usepackage{longtable,booktabs,array}
\usepackage{calc} % for calculating minipage widths
% Correct order of tables after \paragraph or \subparagraph
\usepackage{etoolbox}
\makeatletter
\patchcmd\longtable{\par}{\if@noskipsec\mbox{}\fi\par}{}{}
\makeatother
% Allow footnotes in longtable head/foot
\IfFileExists{footnotehyper.sty}{\usepackage{footnotehyper}}{\usepackage{footnote}}
\makesavenoteenv{longtable}
\setlength{\emergencystretch}{3em} % prevent overfull lines
\providecommand{\tightlist}{%
  \setlength{\itemsep}{0pt}\setlength{\parskip}{0pt}}
\setcounter{secnumdepth}{5}
\usepackage{fontspec}
\setmainfont{DejaVu Serif}
\setsansfont{DejaVu Sans}
\setmonofont[Scale=0.85]{DejaVu Sans Mono}
\usepackage[a4paper,margin=2.2cm]{geometry}
\usepackage{xcolor}
\definecolor{codebg}{HTML}{F4F7FB}
\definecolor{coderule}{HTML}{D3DCE8}
\definecolor{kw}{HTML}{6A3FB0}
\definecolor{str}{HTML}{0B6B50}
\definecolor{com}{HTML}{546880}
\usepackage{minted}
\setminted{breaklines,breakanywhere,fontsize=\small,bgcolor=codebg,frame=single,rulecolor=coderule,framesep=4pt,style=friendly}
\usepackage{titlesec}
\titleformat{\section}{\sffamily\Large\bfseries}{\thesection}{0.6em}{}
\titleformat{\subsection}{\sffamily\large\bfseries}{\thesubsection}{0.6em}{}
\titleformat{\subsubsection}{\sffamily\normalsize\bfseries}{\thesubsubsection}{0.6em}{}
\definecolor{accent}{HTML}{1A58A5}
\AtBeginDocument{\hypersetup{colorlinks=true,linkcolor=accent,urlcolor=accent}}
\usepackage{fancyhdr}
\pagestyle{fancy}
\fancyhf{}
\fancyhead[L]{\sffamily\small Destructuring Field Guide}
\fancyhead[R]{\sffamily\small Şevval Boyraz, 2200004208}
\fancyfoot[C]{\sffamily\small\thepage}
\setlength{\headheight}{14pt}
\sloppy
\ifLuaTeX
  \usepackage{selnolig}  % disable illegal ligatures
\fi
\IfFileExists{bookmark.sty}{\usepackage{bookmark}}{\usepackage{hyperref}}
\IfFileExists{xurl.sty}{\usepackage{xurl}}{} % add URL line breaks if available
\urlstyle{same}
\hypersetup{
  pdftitle={Destructuring Field Guide},
  pdfauthor={Şevval Boyraz, 2200004208},
  hidelinks,
  pdfcreator={LaTeX via pandoc}}

\title{Destructuring Field Guide}
\usepackage{etoolbox}
\makeatletter
\providecommand{\subtitle}[1]{% add subtitle to \maketitle
  \apptocmd{\@title}{\par {\large #1 \par}}{}{}
}
\makeatother
\subtitle{SEN0415 Advanced Web Programming, Assignment 4}
\author{Şevval Boyraz, 2200004208}
\date{20 November 2026}

\begin{document}
\maketitle

{
\setcounter{tocdepth}{3}
\tableofcontents
}
Name: Şevval Boyraz\\
Student number: 2200004208

Bonuses claimed: \textbf{Bonus 1: LaTeX}, \textbf{Bonus 2: Depth},
\textbf{Bonus 3: Real code}, \textbf{Bonus 4: Git repository}

Bonus 4: Git repository ---
\url{https://github.com/sevvalingg/assignment4}

\hypertarget{run}{%
\section{Run}\label{run}}

\begin{minted}{text}
npm install
npm run all      # runs the JS and TS examples
npm run check    # type-checks in strict mode
npm run bonus    # the Bonus 2 experiments (not part of npm run all)
\end{minted}

\texttt{output.\allowbreak{}txt} is the real output of
\texttt{npm run all} on my machine. Two lines that you see in the
terminal are not in \texttt{output.\allowbreak{}txt}: React prints its
warnings (about \texttt{key} and about a function as a child, see P2 and
P3) to stderr, and \texttt{>} only saves stdout.

\hypertarget{my-data-istanbul-ferries}{%
\section{My data: Istanbul ferries}\label{my-data-istanbul-ferries}}

All examples use İstanbul ferry lines (Şehir Hatları): lines are arrays
of piers, a sailing is an object with a nested pier, a timetable is an
object of \texttt{time -> line}, and my trips per pier are a
\texttt{Map}. The piers of the lines are the real ones (I checked them:
the Kadıköy--Karaköy line does not stop at Haydarpaşa any more, and the
Haliç line goes Beşiktaş -- Kabataş -- Karaköy -- Kasımpaşa -- Hasköy --
Sütlüce -- Eyüp). The departure times in the timetable are sample times,
not the real schedule.

\begin{minted}{javascript}
const sailing = {
  line: 'Kadıköy-Karaköy-Eminönü',
  departure: '08:15',
  operator: 'Şehir Hatları',
  pier: { name: 'Kadıköy', side: 'Anadolu' },
  stops: ['Kadıköy', 'Karaköy', 'Eminönü'],
}
\end{minted}

\hypertarget{field-guide}{%
\section{Field guide}\label{field-guide}}

Every entry has the same parts: syntax, example, real output, the same
code without destructuring (one plain assignment per variable, which is
what the pattern really does), \emph{use it when} and \emph{watch out}.

\hypertarget{form-1-values-by-position-const-x-y--f}{%
\subsection{\texorpdfstring{Form 1: values by position,
\texttt{const [\allowbreak{}x,\allowbreak{} y] =\allowbreak{} f(\allowbreak{})}}{Form 1: values by position, }}\label{form-1-values-by-position-const-x-y--f}}

\textbf{Syntax:}
\texttt{const [\allowbreak{}nextTime,\allowbreak{} nextLine] =\allowbreak{} nextFerry(\allowbreak{})}

\textbf{Example} (\texttt{js/\allowbreak{}arrays.\allowbreak{}mjs}):

\begin{minted}{javascript}
function nextFerry() {
  return ['08:15', 'Kadıköy-Karaköy-Eminönü']
}
const [nextTime, nextLine] = nextFerry()
console.log('form 1 (ferry):', nextTime, nextLine)
\end{minted}

\textbf{Output:}
\texttt{form 1 (\allowbreak{}ferry):\allowbreak{} 08:\allowbreak{}15 Kadıköy-Karaköy-Eminönü}

\textbf{Without destructuring:}

\begin{minted}{javascript}
const result = nextFerry()
const nextTime = result[0]
const nextLine = result[1]
\end{minted}

\textbf{Use it when:} a function returns a few values that belong
together and their order is fixed, like \texttt{useState} and
\texttt{useReducer} in React.

\textbf{Watch out:} the names come from the position, not from the data.
If the function changes the order, the names lie and nothing warns you
in JavaScript. In TypeScript, give the function a tuple return type
(\texttt{[\allowbreak{}time:\allowbreak{} string,\allowbreak{} line:\allowbreak{} string]});
without it the return type is \texttt{string[\allowbreak{}]} and every
item is \texttt{string |\allowbreak{} undefined} (see the table).

\hypertarget{form-2-skip-an-item-and-collect-the-rest-const--second-others--list}{%
\subsection{\texorpdfstring{Form 2: skip an item and collect the rest,
\texttt{const [\allowbreak{},\allowbreak{} second,\allowbreak{} .\allowbreak{}.\allowbreak{}.\allowbreak{}others] =\allowbreak{} list}}{Form 2: skip an item and collect the rest, }}\label{form-2-skip-an-item-and-collect-the-rest-const--second-others--list}}

\textbf{Syntax:}
\texttt{const [\allowbreak{},\allowbreak{} firstStop,\allowbreak{} .\allowbreak{}.\allowbreak{}.\allowbreak{}remainingStops] =\allowbreak{} halicLine}

\textbf{Example} (\texttt{js/\allowbreak{}arrays.\allowbreak{}mjs}):

\begin{minted}{javascript}
const halicLine = ['Beşiktaş', 'Kabataş', 'Karaköy', 'Kasımpaşa', 'Hasköy', 'Sütlüce', 'Eyüp']
const [, firstStop, ...remainingStops] = halicLine
console.log('form 2:', firstStop, remainingStops)

const shortLine = ['Kabataş']
const [, onlyStop, ...afterThat] = shortLine
console.log('form 2 (one pier):', onlyStop, afterThat)
\end{minted}

\textbf{Output:}

\begin{minted}{text}
form 2: Kabataş [ 'Karaköy', 'Kasımpaşa', 'Hasköy', 'Sütlüce', 'Eyüp' ]
form 2 (one pier): undefined []
\end{minted}

\textbf{Without destructuring:}

\begin{minted}{javascript}
const firstStop = halicLine[1]
const remainingStops = halicLine.slice(2)
\end{minted}

\textbf{Use it when:} the first item has a different role (here: the
pier the ferry leaves from) and you want "the next one" and "all the
others".

\textbf{Watch out:} the empty place before the first comma is easy to
miss when you read the code. A missing item is not an error, it is
\texttt{undefined}, but the rest is never \texttt{undefined}: it is
always an array, an empty one when nothing is left. The rest must be
last
(\texttt{[\allowbreak{}.\allowbreak{}.\allowbreak{}.\allowbreak{}a,\allowbreak{} b]}
is a syntax error). Rest in a pattern
(\texttt{.\allowbreak{}.\allowbreak{}.\allowbreak{}remainingStops}, on
the left) collects items into a new array; spread in an expression
(\texttt{[\allowbreak{}.\allowbreak{}.\allowbreak{}.\allowbreak{}halicLine]},
on the right) expands an array into items. Same three dots, opposite
jobs.

\hypertarget{form-3-swap-two-variables-a-b--b-a}{%
\subsection{\texorpdfstring{Form 3: swap two variables,
\texttt{[\allowbreak{}a,\allowbreak{} b] =\allowbreak{} [\allowbreak{}b,\allowbreak{} a]}}{Form 3: swap two variables, }}\label{form-3-swap-two-variables-a-b--b-a}}

\textbf{Syntax:}
\texttt{[\allowbreak{}departure,\allowbreak{} arrival] =\allowbreak{} [\allowbreak{}arrival,\allowbreak{} departure]}

\textbf{Example} (\texttt{js/\allowbreak{}arrays.\allowbreak{}mjs}):

\begin{minted}{javascript}
let departure = 'Beşiktaş'
let arrival = 'Eyüp'
console.log('form 3 (going):', departure, '->', arrival)
;[departure, arrival] = [arrival, departure]
console.log('form 3 (return):', departure, '->', arrival)
\end{minted}

\textbf{Output:}

\begin{minted}{text}
form 3 (going): Beşiktaş -> Eyüp
form 3 (return): Eyüp -> Beşiktaş
\end{minted}

\textbf{Without destructuring:}

\begin{minted}{javascript}
const temp = departure
departure = arrival
arrival = temp
\end{minted}

\textbf{Use it when:} two variables must change places (here: the return
trip of the line).

\textbf{Watch out:} the right side is built first, as a new array
\texttt{[\allowbreak{}arrival,\allowbreak{} departure]}, and only then
taken apart, which is why no temporary variable is needed. The variables
must be \texttt{let}. Without semicolons in the file, a line that starts
with \texttt{[\allowbreak{}} is glued to the line before it: that is why
there is a \texttt{;} at the start of the line (see pitfall 6).

\hypertarget{form-4-default-values-in-arrays-const-a--1-b--2--}{%
\subsection{\texorpdfstring{Form 4: default values in arrays,
\texttt{const [\allowbreak{}a =\allowbreak{} 1,\allowbreak{} b =\allowbreak{} 2] =\allowbreak{} [\allowbreak{}]}}{Form 4: default values in arrays, }}\label{form-4-default-values-in-arrays-const-a--1-b--2--}}

\textbf{Syntax:}
\texttt{const [\allowbreak{}from =\allowbreak{} 'Eminönü',\allowbreak{} to =\allowbreak{} 'Kadıköy'] =\allowbreak{} [\allowbreak{}]}

\textbf{Example} (\texttt{js/\allowbreak{}arrays.\allowbreak{}mjs}):

\begin{minted}{javascript}
const [from = 'Eminönü', to = 'Kadıköy'] = []
console.log('form 4 (empty):', from, '->', to)

const [from2 = 'Eminönü', to2 = 'Kadıköy'] = ['Beşiktaş']
console.log('form 4 (one given):', from2, '->', to2)

const [from3 = 'Eminönü', to3 = 'Kadıköy'] = [undefined, null]
console.log('form 4 (undefined, null):', from3, '->', to3)
\end{minted}

\textbf{Output:}

\begin{minted}{text}
form 4 (empty): Eminönü -> Kadıköy
form 4 (one given): Beşiktaş -> Kadıköy
form 4 (undefined, null): Eminönü -> null
\end{minted}

\textbf{Without destructuring:}

\begin{minted}{javascript}
const from = list[0] === undefined ? 'Eminönü' : list[0]
const to = list[1] === undefined ? 'Kadıköy' : list[1]
\end{minted}

\textbf{Use it when:} a list may be shorter than you need and there is a
sensible fallback.

\textbf{Watch out:} the default is used only for \texttt{undefined}.
\texttt{null} is a real value and stays \texttt{null} (last line of the
output). The plain version above shows it: the check is
\texttt{=\allowbreak{}=\allowbreak{}=\allowbreak{} undefined}, not "is
empty".

\hypertarget{form-5-properties-by-name-let--a-b---t}{%
\subsection{\texorpdfstring{Form 5: properties by name,
\texttt{let \{ a,\allowbreak{} b \} =\allowbreak{} t}}{Form 5: properties by name, }}\label{form-5-properties-by-name-let--a-b---t}}

\textbf{Syntax:}
\texttt{const \{ line,\allowbreak{} departure \} =\allowbreak{} sailing}

\textbf{Example} (\texttt{js/\allowbreak{}objects.\allowbreak{}mjs}):

\begin{minted}{javascript}
const { line, departure } = sailing
console.log('form 5 (ferry):', line, departure)
\end{minted}

\textbf{Output:}
\texttt{form 5 (\allowbreak{}ferry):\allowbreak{} Kadıköy-Karaköy-Eminönü 08:\allowbreak{}15}

\textbf{Without destructuring:}

\begin{minted}{javascript}
const line = sailing.line
const departure = sailing.departure
\end{minted}

\textbf{Use it when:} you need a few fields of an object and the field
names are good variable names.

\textbf{Watch out:} the variable name must be the property name; a name
that is not in the object gives \texttt{undefined} without an error (a
typo like \texttt{\{ departue \}} is silent in JavaScript, TypeScript
catches it, see T5). The order inside \texttt{\{ \}} does not matter
(pitfall 5).

\hypertarget{form-6-rename-while-you-take-it-const--a-alpha---t}{%
\subsection{\texorpdfstring{Form 6: rename while you take it,
\texttt{const \{ a:\allowbreak{} alpha \} =\allowbreak{} t}}{Form 6: rename while you take it, }}\label{form-6-rename-while-you-take-it-const--a-alpha---t}}

\textbf{Syntax:}
\texttt{const \{ line:\allowbreak{} eveningLine,\allowbreak{} departure:\allowbreak{} eveningTime \} =\allowbreak{} evening}

\textbf{Example} (\texttt{js/\allowbreak{}objects.\allowbreak{}mjs}):

\begin{minted}{javascript}
// `line` and `departure` already exist from form 5, so I must rename.
const evening = { line: 'Üsküdar-Eminönü', departure: '19:40' }
const { line: eveningLine, departure: eveningTime } = evening
console.log('form 6:', eveningLine, eveningTime)
\end{minted}

\textbf{Output:}
\texttt{form 6:\allowbreak{} Üsküdar-Eminönü 19:\allowbreak{}40}

\textbf{Without destructuring:}

\begin{minted}{javascript}
const eveningLine = evening.line
const eveningTime = evening.departure
\end{minted}

\textbf{Use it when:} the property name is already used by another
variable in the same scope (my case: \texttt{line} exists from form 5,
and \texttt{const \{ line \} =\allowbreak{} evening} would be
\texttt{SyntaxError:\allowbreak{} Identifier 'line' has already been declared}),
or the name is not clear here.

\textbf{Watch out:} read \texttt{line:\allowbreak{} eveningLine} as
"take \texttt{line}, call it \texttt{eveningLine}": the source is on the
left of the colon, the new name on the right. After this, \texttt{line}
from \emph{this} object is not a variable. In TypeScript the colon is
not a type (T1).

\hypertarget{form-7-default-values-in-objects-const--a--1---t}{%
\subsection{\texorpdfstring{Form 7: default values in objects,
\texttt{const \{ a =\allowbreak{} 1 \} =\allowbreak{} t}}{Form 7: default values in objects, }}\label{form-7-default-values-in-objects-const--a--1---t}}

\textbf{Syntax:}
\texttt{const \{ express =\allowbreak{} false,\allowbreak{} capacity =\allowbreak{} 300 \} =\allowbreak{} sailing}

\textbf{Example} (\texttt{js/\allowbreak{}objects.\allowbreak{}mjs}):

\begin{minted}{javascript}
const { express = false, capacity = 300 } = sailing
console.log('form 7 (missing):', express, capacity)

const fastSailing = { line: 'Beşiktaş-Kadıköy', express: true, capacity: undefined }
const { express: isExpress = false, capacity: seats = 300 } = fastSailing
console.log('form 7 (given):', isExpress, seats)
\end{minted}

\textbf{Output:}

\begin{minted}{text}
form 7 (missing): false 300
form 7 (given): true 300
\end{minted}

\textbf{Without destructuring:}

\begin{minted}{javascript}
const express = sailing.express === undefined ? false : sailing.express
const capacity = sailing.capacity === undefined ? 300 : sailing.capacity
\end{minted}

\textbf{Use it when:} a field is optional and most objects do not have
it (most ferries are not express).

\textbf{Watch out:} \texttt{capacity:\allowbreak{} undefined} written in
the object still gets the default, because the rule is "the value is
\texttt{undefined}", not "the key is missing". \texttt{null} does not
get the default (pitfall 2, P6). Rename and default can be combined:
\texttt{capacity:\allowbreak{} seats =\allowbreak{} 300}.

\hypertarget{form-8-rest-in-an-object-const--a-others---t}{%
\subsection{\texorpdfstring{Form 8: rest in an object,
\texttt{const \{ a,\allowbreak{} .\allowbreak{}.\allowbreak{}.\allowbreak{}others \} =\allowbreak{} t}}{Form 8: rest in an object, }}\label{form-8-rest-in-an-object-const--a-others---t}}

\textbf{Syntax:}
\texttt{const \{ operator,\allowbreak{} .\allowbreak{}.\allowbreak{}.\allowbreak{}sailingInfo \} =\allowbreak{} sailing}

\textbf{Example} (\texttt{js/\allowbreak{}objects.\allowbreak{}mjs}):

\begin{minted}{javascript}
const { operator, ...sailingInfo } = sailing
console.log('form 8:', operator)
console.log('form 8 rest:', sailingInfo)
\end{minted}

\textbf{Output:}

\begin{minted}{text}
form 8: Şehir Hatları
form 8 rest: {
  line: 'Kadıköy-Karaköy-Eminönü',
  departure: '08:15',
  pier: { name: 'Kadıköy', side: 'Anadolu' },
  stops: [ 'Kadıköy', 'Karaköy', 'Eminönü' ]
}
\end{minted}

\textbf{Without destructuring:}

\begin{minted}{javascript}
const operator = sailing.operator
const sailingInfo = { ...sailing }
delete sailingInfo.operator
\end{minted}

\textbf{Use it when:} you want to take one or two fields out and pass
"everything else" on. This is exactly what React components do with
\texttt{\{ title,\allowbreak{} children,\allowbreak{} .\allowbreak{}.\allowbreak{}.\allowbreak{}rest \}}
(P4).

\textbf{Watch out:} \texttt{sailingInfo} is a new object, but it is a
shallow copy: \texttt{sailingInfo.\allowbreak{}pier} is the same object
as \texttt{sailing.\allowbreak{}pier} (pitfall 4). The rest must be the
last item in the pattern.

\hypertarget{form-9-nested-patterns-and-arrays-inside-objects-const--p--x--list-first---t}{%
\subsection{\texorpdfstring{Form 9: nested patterns and arrays inside
objects,
\texttt{const \{ p:\allowbreak{} \{ x \},\allowbreak{} list:\allowbreak{} [\allowbreak{}first] \} =\allowbreak{} t}}{Form 9: nested patterns and arrays inside objects, }}\label{form-9-nested-patterns-and-arrays-inside-objects-const--p--x--list-first---t}}

\textbf{Syntax:}
\texttt{const \{ pier:\allowbreak{} \{ name:\allowbreak{} pierName,\allowbreak{} side \},\allowbreak{} stops:\allowbreak{} [\allowbreak{}firstPier,\allowbreak{} ,\allowbreak{} lastPier] \} =\allowbreak{} sailing}

\textbf{Example} (\texttt{js/\allowbreak{}objects.\allowbreak{}mjs}):

\begin{minted}{javascript}
const {
  pier: { name: pierName, side },
  stops: [firstPier, , lastPier],
} = sailing
console.log('form 9:', pierName, side, firstPier, lastPier)

try {
  const nightSailing = { line: 'Kabataş-Kadıköy' }
  const { pier: { name: nightPier } } = nightSailing
  console.log('form 9 (no pier):', nightPier)
} catch (error) {
  console.log('form 9 (no pier):', error.message)
}
\end{minted}

\textbf{Output:}

\begin{minted}{text}
form 9: Kadıköy Anadolu Kadıköy Eminönü
form 9 (no pier): Cannot read properties of undefined (reading 'name')
\end{minted}

\textbf{Without destructuring:}

\begin{minted}{javascript}
const pierName = sailing.pier.name
const side = sailing.pier.side
const firstPier = sailing.stops[0]
const lastPier = sailing.stops[2]
\end{minted}

\textbf{Use it when:} you need a few values from deep inside an object,
and the shape is known.

\textbf{Watch out:} \texttt{pier} and \texttt{stops} are not variables
here; \texttt{pier:\allowbreak{}} only means "go inside". If the inner
object is missing, a missing field (\texttt{undefined}) is fine, but
going \emph{into} it is a \texttt{TypeError} (second output line): the
plain version is \texttt{sailing.\allowbreak{}pier.\allowbreak{}name}
with
\texttt{sailing.\allowbreak{}pier =\allowbreak{}=\allowbreak{}=\allowbreak{} undefined}.
The fix is a default for the inner object:
\texttt{const \{ pier:\allowbreak{} \{ name \} =\allowbreak{} \{\} \} =\allowbreak{} nightSailing}.

\hypertarget{form-10-a-computed-key-const--key-value---t}{%
\subsection{\texorpdfstring{Form 10: a computed key,
\texttt{const \{ [\allowbreak{}key]:\allowbreak{} value \} =\allowbreak{} t}}{Form 10: a computed key, }}\label{form-10-a-computed-key-const--key-value---t}}

\textbf{Syntax:}
\texttt{const \{ [\allowbreak{}wanted]:\allowbreak{} lineAtTime \} =\allowbreak{} timetable}

\textbf{Example} (\texttt{js/\allowbreak{}objects.\allowbreak{}mjs}):

\begin{minted}{javascript}
const timetable = { '07:30': 'Kadıköy-Karaköy-Eminönü', '08:15': 'Üsküdar-Eminönü', '09:00': 'Beşiktaş-Kadıköy' }
const wanted = '08:15'
const { [wanted]: lineAtTime } = timetable
console.log('form 10:', wanted, lineAtTime)

const field = 'operator'
const { [field]: chosen } = sailing
console.log('form 10 (field):', field, chosen)
\end{minted}

\textbf{Output:}

\begin{minted}{text}
form 10: 08:15 Üsküdar-Eminönü
form 10 (field): operator Şehir Hatları
\end{minted}

\textbf{Without destructuring:}

\begin{minted}{javascript}
const lineAtTime = timetable[wanted]
\end{minted}

\textbf{Use it when:} the key is only known at run time (a time the user
picked, a field name from a setting).

\textbf{Watch out:} a computed key always needs a new name after the
colon (\texttt{\{ [\allowbreak{}wanted] \}} alone is a syntax error,
there is no name to give the variable). \texttt{[\allowbreak{}wanted]}
means "the value of \texttt{wanted}", not a key called
\texttt{"wanted"}. Honestly the plain
\texttt{timetable[\allowbreak{}wanted]} is shorter; the pattern is
useful when you take several things at once.

\hypertarget{form-11-patterns-in-parameters-with-defaults-function-f-a--1-b---}{%
\subsection{\texorpdfstring{Form 11: patterns in parameters, with
defaults,
\texttt{function f(\allowbreak{}\{ a =\allowbreak{} 1,\allowbreak{} b \} =\allowbreak{} \{\})}}{Form 11: patterns in parameters, with defaults, }}\label{form-11-patterns-in-parameters-with-defaults-function-f-a--1-b---}}

\textbf{Syntax:}
\texttt{function describeSailing(\allowbreak{}\{ line,\allowbreak{} departure =\allowbreak{} 'saat belli değil',\allowbreak{} express =\allowbreak{} false \} =\allowbreak{} \{\})}

\textbf{Example} (\texttt{js/\allowbreak{}params.\allowbreak{}mjs}):

\begin{minted}{javascript}
function describeSailing({ line, departure = 'saat belli değil', express = false } = {}) {
  const kind = express ? 'ekspres' : 'normal'
  return `${line} (${kind}), kalkış: ${departure}`
}
console.log('form 11 (all):', describeSailing({ line: 'Kadıköy-Beşiktaş', departure: '18:10', express: true }))
console.log('form 11 (only line):', describeSailing({ line: 'Üsküdar-Karaköy-Eminönü' }))
console.log('form 11 (nothing):', describeSailing())

// Form 11b: an array parameter
const trips = [['Kadıköy', 'Eminönü'], ['Beşiktaş', 'Üsküdar'], ['Kabataş', 'Kadıköy']]
const labels = trips.map(([from, to]) => `${from} -> ${to}`)
console.log('form 11b:', labels)
\end{minted}

\textbf{Output:}

\begin{minted}{text}
form 11 (all): Kadıköy-Beşiktaş (ekspres), kalkış: 18:10
form 11 (only line): Üsküdar-Karaköy-Eminönü (normal), kalkış: saat belli değil
form 11 (nothing): undefined (normal), kalkış: saat belli değil
form 11 (no = {}): TypeError - Cannot destructure property 'line' of 'undefined' as it is undefined.
form 11b: [ 'Kadıköy -> Eminönü', 'Beşiktaş -> Üsküdar', 'Kabataş -> Kadıköy' ]
\end{minted}

\textbf{Without destructuring:}

\begin{minted}{javascript}
function describeSailing(options) {
  if (options === undefined) options = {}
  const line = options.line
  const departure = options.departure === undefined ? 'saat belli değil' : options.departure
  const express = options.express === undefined ? false : options.express
  // ...
}
\end{minted}

\textbf{Use it when:} a function takes an "options object": the caller
writes names instead of remembering an order, and can leave things out.
Every React component is this form (P1).

\textbf{Watch out:} there are two defaults: the ones inside the pattern
(for missing fields) and \texttt{=\allowbreak{} \{\}} after it (for a
missing argument). Without \texttt{=\allowbreak{} \{\}}, calling with no
argument throws the \texttt{TypeError} in the fourth output line,
because there is no object to take apart.

\hypertarget{form-12-patterns-in-loops-for-const-k-v-of-objectentrieso-and-map-entries}{%
\subsection{\texorpdfstring{Form 12: patterns in loops,
\texttt{for (\allowbreak{}const [\allowbreak{}k,\allowbreak{} v] of Object.\allowbreak{}entries(\allowbreak{}o))}
and Map
entries}{Form 12: patterns in loops,  and Map entries}}\label{form-12-patterns-in-loops-for-const-k-v-of-objectentrieso-and-map-entries}}

\textbf{Syntax:}
\texttt{for (\allowbreak{}const [\allowbreak{}pier,\allowbreak{} count] of myTrips)}

\textbf{Example} (\texttt{js/\allowbreak{}params.\allowbreak{}mjs}):

\begin{minted}{javascript}
for (const [time, line] of Object.entries(timetable)) {
  console.log('form 12 (object):', time, line)
}

const myTrips = new Map([['Kadıköy', 5], ['Beşiktaş', 2], ['Üsküdar', 1]])
for (const [pier, count] of myTrips) {
  console.log('form 12 (Map):', pier, count)
}
\end{minted}

\textbf{Output:}

\begin{minted}{text}
form 12 (object): 07:30 Kadıköy-Karaköy-Eminönü
form 12 (object): 08:15 Üsküdar-Eminönü
form 12 (object): 09:00 Beşiktaş-Kadıköy
form 12 (Map): Kadıköy 5
form 12 (Map): Beşiktaş 2
form 12 (Map): Üsküdar 1
\end{minted}

\textbf{Without destructuring:}

\begin{minted}{javascript}
for (const entry of myTrips) {
  console.log(entry[0], entry[1])
}
\end{minted}

\textbf{Use it when:} you loop over key--value pairs. Iterating a
\texttt{Map} gives \texttt{[\allowbreak{}key,\allowbreak{} value]}
arrays, and \texttt{Object.\allowbreak{}entries} turns an object into
the same shape.

\textbf{Watch out:} a plain object is not iterable:
\texttt{for (\allowbreak{}const [\allowbreak{}k,\allowbreak{} v] of timetable)}
is a \texttt{TypeError}, you need \texttt{Object.\allowbreak{}entries}.
Each pair is an array, so the pattern uses \texttt{[\allowbreak{} ]},
not \texttt{\{ \}}. \texttt{Object.\allowbreak{}entries} gives keys as
strings, even for number keys.

\hypertarget{pitfalls}{%
\section{Pitfalls}\label{pitfalls}}

All from \texttt{js/\allowbreak{}pitfalls.\allowbreak{}mjs}.

\hypertarget{pitfall-1-destructuring-undefined-or-null-throws}{%
\subsection{\texorpdfstring{Pitfall 1: destructuring \texttt{undefined}
or \texttt{null}
throws}{Pitfall 1: destructuring  or  throws}}\label{pitfall-1-destructuring-undefined-or-null-throws}}

\begin{minted}{javascript}
const { a } = undefined
\end{minted}

\textbf{Output:} \texttt{pitfall 1:\allowbreak{} TypeError}

\textbf{Why:} a pattern reads properties of the value on the right.
\texttt{undefined} and \texttt{null} have no properties at all, so even
the first read (\texttt{undefined.\allowbreak{}a}) fails. Fix: a default
for the whole value (\texttt{=\allowbreak{} \{\}}) or a check before.

\hypertarget{pitfall-2-a-default-only-applies-to-undefined-not-to-null}{%
\subsection{\texorpdfstring{Pitfall 2: a default only applies to
\texttt{undefined}, not to
\texttt{null}}{Pitfall 2: a default only applies to , not to }}\label{pitfall-2-a-default-only-applies-to-undefined-not-to-null}}

\begin{minted}{javascript}
const unknownFerry = { line: 'Kabataş-Kadıköy', capacity: null }
const { capacity = 300 } = unknownFerry
const safeCapacity = unknownFerry.capacity ?? 300
\end{minted}

\textbf{Output:}

\begin{minted}{text}
pitfall 2: null
pitfall 2 (fix with ??): 300
\end{minted}

\textbf{Why:} the language checks
\texttt{value =\allowbreak{}=\allowbreak{}=\allowbreak{} undefined}
before it uses a default. \texttt{null} means "there is a value and it
is nothing", so it is kept. Data from a server often has \texttt{null}
for "unknown"; use \texttt{??} (which treats \texttt{null} and
\texttt{undefined} the same) when \texttt{null} should also get the
fallback.

\hypertarget{pitfall-3-assignment-without-constlet-needs-parentheses}{%
\subsection{\texorpdfstring{Pitfall 3: assignment without
\texttt{const}/\texttt{let} needs
parentheses}{Pitfall 3: assignment without / needs parentheses}}\label{pitfall-3-assignment-without-constlet-needs-parentheses}}

\begin{minted}{javascript}
new Function("let line; { line } = { line: 'Kadıköy-Beşiktaş' }")
// fix:
;({ line } = { line: 'Kadıköy-Beşiktaş' })
\end{minted}

\textbf{Output:}

\begin{minted}{text}
pitfall 3: SyntaxError - Unexpected token '='
pitfall 3 (fix): Kadıköy-Beşiktaş
\end{minted}

\textbf{Why:} at the start of a statement \texttt{\{} opens a block, not
an object pattern, so \texttt{\{ line \}} is a block with the word
\texttt{line} in it, and the \texttt{=\allowbreak{}} after a block makes
no sense. Parentheses make it an expression. (I compile the bad line
with \texttt{new Function} so that the rest of the file still runs;
written directly, the whole file would not start.)

\hypertarget{pitfall-4-it-copies-references-not-objects}{%
\subsection{Pitfall 4: it copies references, not
objects}\label{pitfall-4-it-copies-references-not-objects}}

\begin{minted}{javascript}
const sailing = { line: 'Kadıköy-Karaköy-Eminönü', pier: { name: 'Kadıköy', side: 'Anadolu' } }
const { pier } = sailing
const { line: copiedLine, ...rest } = sailing
pier.name = 'Moda'
\end{minted}

\textbf{Output:}

\begin{minted}{text}
pitfall 4 (original): Moda
pitfall 4 (rest): Moda
pitfall 4 (same object?): true true
\end{minted}

\textbf{Why:} \texttt{const \{ pier \} =\allowbreak{} sailing} is
\texttt{const pier =\allowbreak{} sailing.\allowbreak{}pier}: the
variable points to the \textbf{same} object. Rest makes a new outer
object but copies the inner ones by reference too. To change a copy
safely you need a real copy
(\texttt{structuredClone(\allowbreak{}sailing)} or
\texttt{\{ .\allowbreak{}.\allowbreak{}.\allowbreak{}sailing,\allowbreak{} pier:\allowbreak{} \{ .\allowbreak{}.\allowbreak{}.\allowbreak{}sailing.\allowbreak{}pier \} \}}).

\hypertarget{pitfall-5-in-objects-the-order-does-not-matter-in-arrays-it-does}{%
\subsection{Pitfall 5: in objects the order does not matter, in arrays
it
does}\label{pitfall-5-in-objects-the-order-does-not-matter-in-arrays-it-does}}

\begin{minted}{javascript}
const route = { from: 'Beşiktaş', to: 'Üsküdar' }
const { to, from } = route
const routePair = ['Beşiktaş', 'Üsküdar']
const [arrival, departure] = routePair
\end{minted}

\textbf{Output:}

\begin{minted}{text}
pitfall 5 (object): Beşiktaş -> Üsküdar
pitfall 5 (array): departure = Üsküdar , arrival = Beşiktaş
\end{minted}

\textbf{Why:} object patterns match by name, array patterns by position.
In the array version my names are wrong (Beşiktaş is the departure) and
JavaScript cannot know it. This is the reason React\textquotesingle s
\texttt{useState} returns an array (you choose the names) and most other
APIs return objects (the names are fixed and safe).

\hypertarget{pitfall-6-a-swap-without-a-semicolon-on-the-line-before}{%
\subsection{Pitfall 6: a swap without a semicolon on the line
before}\label{pitfall-6-a-swap-without-a-semicolon-on-the-line-before}}

\begin{minted}{javascript}
let a = 'Eminönü'
let b = 'Kadıköy'
console.log('pitfall 6 before:', a, b)
[a, b] = [b, a]
\end{minted}

\textbf{Output:}

\begin{minted}{text}
pitfall 6 before: Eminönü Kadıköy
pitfall 6: TypeError - Cannot set properties of undefined (setting 'Kadıköy')
\end{minted}

\textbf{Why:} JavaScript only inserts a semicolon automatically when the
next line cannot continue the statement. A \texttt{[\allowbreak{}}
\emph{can} continue it, so the two lines are read as one:
\texttt{console.\allowbreak{}log(\allowbreak{}.\allowbreak{}.\allowbreak{}.\allowbreak{})[\allowbreak{}a,\allowbreak{} b] =\allowbreak{} [\allowbreak{}b,\allowbreak{} a]}.
\texttt{console.\allowbreak{}log} returns \texttt{undefined},
\texttt{[\allowbreak{}a,\allowbreak{} b]} is the comma operator (=
\texttt{b} = \texttt{'Kadıköy'}), and the code tries
\texttt{undefined[\allowbreak{}'Kadıköy'] =\allowbreak{} .\allowbreak{}.\allowbreak{}.\allowbreak{}}.
Fix: start the line with \texttt{;} (as in form 3).

\hypertarget{props-and-children}{%
\section{Props and children}\label{props-and-children}}

Code: \texttt{js/\allowbreak{}props.\allowbreak{}jsx} (P1--P6 and one
extra) and \texttt{ts/\allowbreak{}props.\allowbreak{}tsx} (P1--P7,
typed). Rendered with \texttt{renderToStaticMarkup}, so the output is
the real HTML.

\hypertarget{p1-props-is-one-object}{%
\subsection{P1: props is one object}\label{p1-props-is-one-object}}

\begin{minted}{jsx}
function LineCardA(props) {
  return <section><h2>{props.line}</h2><p>{props.operator}</p>{props.children}</section>
}
function LineCardB({ line, operator, children }) {
  return <section><h2>{line}</h2><p>{operator}</p>{children}</section>
}
function ShowProps(props) {
  console.log('P1 props object:', props)
  return null
}
renderToStaticMarkup(<ShowProps line="Kabataş-Kadıköy" express>sabah seferi</ShowProps>)
\end{minted}

\textbf{Output:}

\begin{minted}{text}
P1 A: <section><h2>Kadıköy-Beşiktaş</h2><p>Şehir Hatları</p>Her 30 dakikada bir</section>
P1 B: <section><h2>Kadıköy-Beşiktaş</h2><p>Şehir Hatları</p>Her 30 dakikada bir</section>
P1 props object: { line: 'Kabataş-Kadıköy', express: true, children: 'sabah seferi' }
\end{minted}

React calls the component as a function with one argument: an object
made from the JSX attributes, plus \texttt{children}. \texttt{LineCardB}
is form 11 (a pattern in the parameter list) applied to that object with
form 5 (properties by name). \texttt{LineCardA} is the same code without
destructuring. The printed object shows that \texttt{children} is simply
one more field of it, next to \texttt{line}, and that \texttt{express}
with no value becomes \texttt{true}.

\hypertarget{p2-children-is-only-a-prop}{%
\subsection{P2: children is only a
prop}\label{p2-children-is-only-a-prop}}

\begin{minted}{jsx}
<LineCardB line="Üsküdar-Eminönü">İlk sefer 07:00</LineCardB>
<LineCardB line="Üsküdar-Eminönü" children="İlk sefer 07:00" />
<LineCardB line="Üsküdar-Eminönü" children="attribute">between tags</LineCardB>
\end{minted}

\textbf{Output:}

\begin{minted}{text}
P2 between tags: <section><h2>Üsküdar-Eminönü</h2><p></p>İlk sefer 07:00</section>
P2 as attribute: <section><h2>Üsküdar-Eminönü</h2><p></p>İlk sefer 07:00</section>
P2 both: <section><h2>Üsküdar-Eminönü</h2><p></p>between tags</section>
\end{minted}

The first two are identical: content between the tags is only a nicer
way to write the \texttt{children} prop. When both are given, the
content between the tags wins: JSX turns into
\texttt{jsx(\allowbreak{}LineCardB,\allowbreak{} \{ line,\allowbreak{} children:\allowbreak{} 'attribute',\allowbreak{} children:\allowbreak{} 'between tags' \})},
and like in any object the later key overwrites the earlier one. (The
empty \texttt{<p></\allowbreak{}p>} is because I gave no
\texttt{operator}.)

\texttt{key} is the opposite case: it looks like a prop but is not one.
React takes it out before the component is called:

\begin{minted}{text}
P2 key: <span>key is undefined</span>
\end{minted}

and prints this warning on stderr (so it is not in
\texttt{output.\allowbreak{}txt}):
\texttt{ShowKey:\allowbreak{} `key` is not a prop.\allowbreak{} Trying to access it will result in `undefined` being returned.\allowbreak{}}

\hypertarget{p3-what-children-can-be}{%
\subsection{P3: what children can be}\label{p3-what-children-can-be}}

Each value rendered inside
\texttt{<div>[\allowbreak{}\{children\}]</\allowbreak{}div>}:

\begin{longtable}[]{@{}>{\raggedright\arraybackslash}p{(\linewidth - 4\tabcolsep) * \real{0.36}}>{\raggedright\arraybackslash}p{(\linewidth - 4\tabcolsep) * \real{0.46}}>{\raggedright\arraybackslash}p{(\linewidth - 4\tabcolsep) * \real{0.18}}@{}}
\toprule\noalign{}
children & output & appears? \\
\midrule\noalign{}
\endhead
\bottomrule\noalign{}
\endlastfoot
\texttt{'Kadıköy'} &
\texttt{<div>[\allowbreak{}Kadıköy]</\allowbreak{}div>} & yes \\
\texttt{42} & \texttt{<div>[\allowbreak{}42]</\allowbreak{}div>} &
yes \\
\texttt{0} & \texttt{<div>[\allowbreak{}0]</\allowbreak{}div>} & yes \\
\texttt{<b>Eminönü</\allowbreak{}b>} &
\texttt{<div>[\allowbreak{}<b>Eminönü</\allowbreak{}b>]</\allowbreak{}div>}
& yes \\
\texttt{[\allowbreak{}<b>Beşiktaş</\allowbreak{}b>,\allowbreak{} <i>Üsküdar</\allowbreak{}i>]}
&
\texttt{<div>[\allowbreak{}<b>Beşiktaş</\allowbreak{}b><i>Üsküdar</\allowbreak{}i>]</\allowbreak{}div>}
& yes, one after the other \\
\texttt{undefined} & \texttt{<div>[\allowbreak{}]</\allowbreak{}div>} &
no \\
\texttt{null} & \texttt{<div>[\allowbreak{}]</\allowbreak{}div>} & no \\
\texttt{false} & \texttt{<div>[\allowbreak{}]</\allowbreak{}div>} &
no \\
\texttt{true} & \texttt{<div>[\allowbreak{}]</\allowbreak{}div>} & no \\
\texttt{''} & \texttt{<div>[\allowbreak{}]</\allowbreak{}div>} & no
(empty text) \\
\texttt{(\allowbreak{}) =\allowbreak{}> 'Kabataş'} &
\texttt{<div>[\allowbreak{}]</\allowbreak{}div>} & no, and a warning \\
\end{longtable}

\texttt{undefined}, \texttt{null} and booleans render nothing on
purpose, so that \texttt{\{isExpress \&\& <b>ekspres</\allowbreak{}b>\}}
works: when the condition is \texttt{false} nothing appears. The
surprise is \texttt{0}: it is a number, so it is printed.
\texttt{\{count \&\& <b>.\allowbreak{}.\allowbreak{}.\allowbreak{}</\allowbreak{}b>\}}
with \texttt{count =\allowbreak{} 0} shows a \texttt{0} on the screen;
write
\texttt{\{count > 0 \&\& .\allowbreak{}.\allowbreak{}.\allowbreak{}\}}.
A function is not something React can render; it warns
\texttt{Functions are not valid as a React child} and renders nothing. A
function \emph{does} work as children when the component calls it (a
"render prop"):

\begin{minted}{jsx}
function NextSailing({ children }) {
  const time = '08:15'
  return <p>{children(time)}</p>
}
<NextSailing>{(time) => `Sonraki vapur ${time}`}</NextSailing>
\end{minted}

\textbf{Output:}
\texttt{P3 function,\allowbreak{} called:\allowbreak{} <p>Sonraki vapur 08:\allowbreak{}15</\allowbreak{}p>}

\hypertarget{p4-separating-children-from-the-other-props}{%
\subsection{P4: separating children from the other
props}\label{p4-separating-children-from-the-other-props}}

\begin{minted}{jsx}
function PierSection({ title, children, ...rest }) {
  return <section {...rest}><h3>{title}</h3>{children}</section>
}
<PierSection title="Kadıköy İskelesi" id="kadikoy" className="pier" data-side="Anadolu">Beşiktaş, Eminönü, Karaköy</PierSection>

function SpreadThenExplicit({ children, ...rest }) {
  return <section {...rest} className="fixed">{children}</section>
}
function ExplicitThenSpread({ children, ...rest }) {
  return <section className="fixed" {...rest}>{children}</section>
}
\end{minted}

\textbf{Output:}

\begin{minted}{text}
P4 rest: <section id="kadikoy" class="pier" data-side="Anadolu"><h3>Kadıköy İskelesi</h3>Beşiktaş, Eminönü, Karaköy</section>
P4 spread then explicit: <section class="fixed">x</section>
P4 explicit then spread: <section class="from-parent">x</section>
P4 title not in rest: <section lang="tr"><h3>Beşiktaş</h3>x</section>
\end{minted}

\texttt{\{ title,\allowbreak{} children,\allowbreak{} .\allowbreak{}.\allowbreak{}.\allowbreak{}rest \}}
is form 8 on the props object: the two props the component uses itself
are taken out by name, and every other prop (\texttt{id},
\texttt{className}, \texttt{data-side}) goes into \texttt{rest} and is
spread onto the real \texttt{<section>}. This is how a component lets
the caller add any HTML attribute without listing them all.

The order of \texttt{\{.\allowbreak{}.\allowbreak{}.\allowbreak{}rest\}}
and an explicit prop decides who wins: the last one written wins,
exactly like keys in an object literal
(\texttt{\{ .\allowbreak{}.\allowbreak{}.\allowbreak{}rest,\allowbreak{} className:\allowbreak{} 'fixed' \}}
vs
\texttt{\{ className:\allowbreak{} 'fixed',\allowbreak{} .\allowbreak{}.\allowbreak{}.\allowbreak{}rest \}}).
Put the explicit prop after the spread when the component must keep
control, before it when the caller may override it. A name taken out by
the pattern (\texttt{title}) is not in \texttt{rest}, so it does not
leak onto the \texttt{<section>} as an HTML attribute.

\hypertarget{p5-props-or-children}{%
\subsection{P5: props or children?}\label{p5-props-or-children}}

\begin{minted}{jsx}
function TimetableWithProps({ header, footer, children }) {
  return (
    <article>
      <header>{header}</header>
      <ul>{children}</ul>
      <footer>{footer}</footer>
    </article>
  )
}
function TimetableWithChildren({ children }) {
  return <article>{children}</article>
}
\end{minted}

\textbf{Output:}

\begin{minted}{text}
P5 props: <article><header>Kadıköy-Karaköy-Eminönü</header><ul><li>07:30</li><li>08:15</li></ul><footer>Tarife değişebilir</footer></article>
P5 children: <article><header>Kadıköy-Karaköy-Eminönü</header><ul><li>07:30</li><li>08:15</li></ul><footer>Tarife değişebilir</footer></article>
\end{minted}

The HTML is the same; the difference is who decides the structure. With
\texttt{header} and \texttt{footer} props, the component owns the
layout: every timetable card has the header first, the list in a
\texttt{<ul>}, the footer last, and the caller cannot get it wrong. With
children only, the component is just a frame and the caller writes
everything inside, so the caller can also forget the footer or put
things in another order. I would use props for named parts that every
card must have in fixed places (a header, a footer, an icon) and
children for the main free content, which is what
\texttt{TimetableWithProps} does: named slots as props, the list as
children.

\hypertarget{p6-defaults-for-props}{%
\subsection{P6: defaults for props}\label{p6-defaults-for-props}}

\begin{minted}{jsx}
function SailingBadge({ kind = 'normal sefer' }) {
  return <span className="badge">{kind}</span>
}
<SailingBadge />
<SailingBadge kind={undefined} />
<SailingBadge kind={null} />
<SailingBadge kind="ekspres" />
\end{minted}

\textbf{Output:}

\begin{minted}{text}
P6 missing: <span class="badge">normal sefer</span>
P6 undefined: <span class="badge">normal sefer</span>
P6 null: <span class="badge"></span>
P6 given: <span class="badge">ekspres</span>
\end{minted}

This is form 7 on the props object. A missing prop and
\texttt{kind=\allowbreak{}\{undefined\}} are the same thing (the value
is \texttt{undefined}), so both get the default.
\texttt{kind=\allowbreak{}\{null\}} is a value, the default is not used,
and \texttt{null} renders nothing (P3): on the screen the user sees an
empty badge, a coloured box with no text. If \texttt{null} can come from
data, use \texttt{kind ?? 'normal sefer'} in the body.

\hypertarget{extra-a-component-that-changes-its-props}{%
\subsection{Extra: a component that changes its
props}\label{extra-a-component-that-changes-its-props}}

\begin{minted}{jsx}
function ChangesProps(props) {
  props.line = 'changed'
  // ...
}
\end{minted}

\textbf{Output:}
\texttt{extra change props:\allowbreak{} <p>TypeError:\allowbreak{} Cannot assign to read only property \&\#x27;line\&\#x27; of object \&\#x27;\#\&lt;Object\&gt;\&\#x27;</\allowbreak{}p>}

In development React freezes the props object, so a component cannot
change what it received. Props are input; if a value must change, it
belongs in state.

\hypertarget{p7-types-for-props-and-children-tspropstsx}{%
\subsection{\texorpdfstring{P7: types for props and children
(\texttt{ts/\allowbreak{}props.\allowbreak{}tsx})}{P7: types for props and children ()}}\label{p7-types-for-props-and-children-tspropstsx}}

\textbf{(a) \texttt{children?:\allowbreak{} ReactNode}, written by
hand.} \texttt{ReactNode} is "anything React can render": string,
number, element, array of those, \texttt{null}, \texttt{undefined},
boolean. With \texttt{?} the component can also be used with nothing
inside.

\begin{minted}{tsx}
type NoteProps = { pier: string; children?: ReactNode }
function Note({ pier, children }: NoteProps) {
  return <p><b>{pier}:</b> {children}</p>
}
\end{minted}

Output:
\texttt{P7 ReactNode:\allowbreak{} <p><b>Beşiktaş:\allowbreak{}</\allowbreak{}b> iskele yenileniyor</\allowbreak{}p>}
and
\texttt{P7 ReactNode,\allowbreak{} empty:\allowbreak{} <p><b>Beşiktaş:\allowbreak{}</\allowbreak{}b> </\allowbreak{}p>}.
Rejected:
\texttt{<Note pier=\allowbreak{}"Beşiktaş">\{(\allowbreak{}) =\allowbreak{}> 'x'\}</\allowbreak{}Note>}
gives
\texttt{error TS2322:\allowbreak{} Type '(\allowbreak{}) =\allowbreak{}> string' is not assignable to type 'ReactNode'.\allowbreak{}}
That is P3\textquotesingle s function case caught before running. If you
want a function as children, the type must say so:
\texttt{children:\allowbreak{} (\allowbreak{}time:\allowbreak{} string) =\allowbreak{}> ReactNode}
(the \texttt{NextSailing} render prop in
\texttt{ts/\allowbreak{}props.\allowbreak{}tsx}).

\textbf{(b) \texttt{PropsWithChildren<P>}} adds
\texttt{children?:\allowbreak{} ReactNode} to your own props type.

\begin{minted}{tsx}
type CrossingProps = PropsWithChildren<{ from: string; to: string }>
\end{minted}

Output:
\texttt{P7 PropsWithChildren:\allowbreak{} <p>Kabataş → Kadıköy (\allowbreak{}20 dk)</\allowbreak{}p>}.
Rejected:
\texttt{<Crossing from=\allowbreak{}"Kabataş">x</\allowbreak{}Crossing>}
gives
\texttt{error TS2741:\allowbreak{} Property 'to' is missing in type '\{ children:\allowbreak{} string; from:\allowbreak{} string; \}' but required in type '\{ from:\allowbreak{} string; to:\allowbreak{} string; \}'.\allowbreak{}}
Note how the compiler shows \texttt{children} as one of the props: the
type checker agrees with P2.

\textbf{(c) \texttt{ComponentProps<'button'>}} is every prop a real
\texttt{<button>} accepts (\texttt{disabled}, \texttt{onClick},
\texttt{aria-label}, \texttt{children}, ...), so a wrapper can accept
them all and pass them on with rest/spread.

\begin{minted}{tsx}
type TicketButtonProps = ComponentProps<'button'> & { price: number }
function TicketButton({ price, children, type = 'button', ...rest }: TicketButtonProps) {
  return <button type={type} {...rest}>{children} ({price} TL)</button>
}
\end{minted}

Output:
\texttt{P7 ComponentProps:\allowbreak{} <button type=\allowbreak{}"button" disabled=\allowbreak{}"" aria-label=\allowbreak{}"bilet al">Bilet al (\allowbreak{}59 TL)</\allowbreak{}button>}.
Rejected:
\texttt{<TicketButton price=\allowbreak{}\{59\} href=\allowbreak{}"/\allowbreak{}bilet">}
gives
\texttt{error TS2322:\allowbreak{} .\allowbreak{}.\allowbreak{}.\allowbreak{} Property 'href' does not exist on type '.\allowbreak{}.\allowbreak{}.\allowbreak{}'.\allowbreak{} Did you mean 'ref'?}
(a button has no \texttt{href}), and \texttt{<TicketButton>} without
\texttt{price} gives
\texttt{error TS2741:\allowbreak{} Property 'price' is missing in type '\{ children:\allowbreak{} string; \}' but required in type '\{ price:\allowbreak{} number; \}'.\allowbreak{}}

The rejected lines are written as comments in
\texttt{ts/\allowbreak{}props.\allowbreak{}tsx} with these messages,
because the handout allows \texttt{@ts-expect-error} only in T5. I got
the messages by compiling each line once.

Summary of the difference: props are the named inputs of a component
(attributes in JSX); \textbf{children} is one specific prop that holds
what you write between the tags. Use named props when the component must
know \emph{what} each thing is and \emph{where} it goes; use children
when the component only wraps content it does not need to understand.

\hypertarget{typescript-topics}{%
\section{TypeScript topics}\label{typescript-topics}}

Code: \texttt{ts/\allowbreak{}types.\allowbreak{}ts}.

\hypertarget{t1-in-a-pattern--a-string--is-a-rename-not-a-type}{%
\subsection{\texorpdfstring{T1: in a pattern,
\texttt{\{ a:\allowbreak{} string \}} is a rename, not a
type}{T1: in a pattern,  is a rename, not a type}}\label{t1-in-a-pattern--a-string--is-a-rename-not-a-type}}

\begin{minted}{typescript}
const sailing = { line: 'Kadıköy-Karaköy-Eminönü', departure: '08:15', capacity: 300 }
const { line: string } = sailing
console.log('T1 (confusion): a variable named string =', string)

const { departure, capacity }: { departure: string; capacity: number } = sailing
\end{minted}

\textbf{Output:}

\begin{minted}{text}
T1 (confusion): a variable named string = Kadıköy-Karaköy-Eminönü
T1 (correct): 08:15 300
\end{minted}

The first line compiles, which is the dangerous part: it creates a
variable called \texttt{string}. Inside a pattern the colon always means
"rename" (form 6). The type is written once, after the whole pattern,
describing the whole object. With two fields renamed to \texttt{string}
you would get
\texttt{error TS2451:\allowbreak{} Cannot redeclare block-scoped variable 'string'.\allowbreak{}}

\hypertarget{t2-tuple-return-types-and-as-const}{%
\subsection{\texorpdfstring{T2: tuple return types and
\texttt{as const}}{T2: tuple return types and }}\label{t2-tuple-return-types-and-as-const}}

\begin{minted}{typescript}
function nextSailing(): [time: string, line: string] {
  return ['08:15', 'Üsküdar-Eminönü']
}
const [nextTime, nextLine] = nextSailing()       // string, string

const plainPair = ['Kadıköy', 'Eminönü']          // string[]
const [plainFrom] = plainPair                      // string | undefined

const fixedPair = ['Kadıköy', 'Eminönü'] as const // readonly ['Kadıköy', 'Eminönü']
const [fixedFrom, fixedTo] = fixedPair             // 'Kadıköy', 'Eminönü'
const homePier: 'Kadıköy' = fixedFrom
\end{minted}

\textbf{Output:}

\begin{minted}{text}
T2 (tuple): 08:15 Üsküdar-Eminönü
T2 (no as const): Kadıköy
T2 (as const): Kadıköy Eminönü
\end{minted}

A tuple type fixes the length and the type of each position, so
destructured items are never \texttt{undefined}. An array literal
without annotation is just \texttt{string[\allowbreak{}]}.
\texttt{as const} makes it a readonly tuple of literal types:
\texttt{fixedFrom} has the type \texttt{'Kadıköy'}, which is why it can
be assigned to a variable of type \texttt{'Kadıköy'}; \texttt{plainFrom}
could not.

\hypertarget{t3-nouncheckedindexedaccess}{%
\subsection{\texorpdfstring{T3:
\texttt{noUncheckedIndexedAccess}}{T3: }}\label{t3-nouncheckedindexedaccess}}

\begin{minted}{typescript}
const todaysPiers: string[] = ['Beşiktaş', 'Üsküdar']
const [firstPier] = todaysPiers                    // string | undefined
if (firstPier !== undefined) {
  console.log('T3 (check):', firstPier.toUpperCase()) // narrowed to string here
}
console.log('T3 (??):', (firstPier ?? 'yok').toUpperCase())
const [pierWithDefault = 'yok'] = todaysPiers      // string
\end{minted}

\textbf{Output:}
\texttt{T3 (\allowbreak{}check):\allowbreak{} BEŞIKTAŞ},
\texttt{T3 (\allowbreak{}??):\allowbreak{} BEŞIKTAŞ},
\texttt{T3 (\allowbreak{}default):\allowbreak{} BEŞIKTAŞ}

With this option on, every item read from an array of unknown length
gets \texttt{|\allowbreak{} undefined}, because the array could be
shorter (form 2 showed that at run time). Calling a method on it
directly is
\texttt{error TS18048:\allowbreak{} 'firstPier' is possibly 'undefined'.\allowbreak{}}
(I got this exact error when I tried
\texttt{firstStop.\allowbreak{}toUpperCase(\allowbreak{})} in
\texttt{ts/\allowbreak{}arrays.\allowbreak{}ts}.) Three fixes: check it
(TypeScript narrows the type inside the \texttt{if}), use \texttt{??},
or give the pattern a default.

\hypertarget{t4-optional-fields-and-defaults-change-the-type}{%
\subsection{T4: optional fields and defaults change the
type}\label{t4-optional-fields-and-defaults-change-the-type}}

\begin{minted}{typescript}
type FerryOptions = { express?: boolean; note?: string | null }
const options: FerryOptions = { note: null }
const { express } = options                        // boolean | undefined
const { express: expressOrFalse = false } = options // boolean
const { note = 'not yok' } = options               // string | null
\end{minted}

\textbf{Output:} \texttt{T4:\allowbreak{} undefined false null}

An optional field (\texttt{?}) adds \texttt{undefined} to the type of
the variable. A default removes \texttt{undefined}, because TypeScript
knows the run-time rule (form 7). It does not remove \texttt{null},
because at run time \texttt{null} does not trigger the default (pitfall
2): the type and the behaviour agree, and the output shows
\texttt{null}.

\hypertarget{t5-one-compile-error-caused-by-destructuring}{%
\subsection{T5: one compile error caused by
destructuring}\label{t5-one-compile-error-caused-by-destructuring}}

\begin{minted}{typescript}
// @ts-expect-error -- `platform` is not in the type of `sailing`
const { platform } = sailing
\end{minted}

Without the comment line:

\begin{minted}{text}
error TS2339: Property 'platform' does not exist on type '{ line: string; departure: string; capacity: number; }'.
\end{minted}

\textbf{Output with the comment:}
\texttt{T5 (\allowbreak{}runs anyway,\allowbreak{} but):\allowbreak{} undefined}

In JavaScript a name that is not in the object silently gives
\texttt{undefined}; TypeScript refuses it, which catches typos.
\texttt{@ts-expect-error} silences exactly one error on the next line
(and is itself an error if there is none). The fix is to put the field
in the type, as optional because not every sailing has one:

\begin{minted}{typescript}
type SailingWithPlatform = { line: string; departure: string; platform?: number }
const { platform: platformNumber = 1 } = withPlatform
\end{minted}

\textbf{Output:} \texttt{T5 (\allowbreak{}fix):\allowbreak{} 2}

\hypertarget{javascript-vs-typescript}{%
\section{JavaScript vs TypeScript}\label{javascript-vs-typescript}}

What stays the same, what TypeScript needs in addition, and the real
error when it is missing or wrong. Every error below I produced by
compiling the wrong line with \texttt{npm run check}.

\begin{longtable}[]{@{}>{\raggedright\arraybackslash}p{(\linewidth - 6\tabcolsep) * \real{0.13}}>{\raggedright\arraybackslash}p{(\linewidth - 6\tabcolsep) * \real{0.17}}>{\raggedright\arraybackslash}p{(\linewidth - 6\tabcolsep) * \real{0.30}}>{\raggedright\arraybackslash}p{(\linewidth - 6\tabcolsep) * \real{0.40}}@{}}
\toprule\noalign{}
Form & Same in TypeScript & What you add & Error when you get it
wrong \\
\midrule\noalign{}
\endhead
\bottomrule\noalign{}
\endlastfoot
1
\texttt{const [\allowbreak{}a,\allowbreak{} b] =\allowbreak{} f(\allowbreak{})}
& the pattern & a tuple return type
\texttt{[\allowbreak{}time:\allowbreak{} string,\allowbreak{} line:\allowbreak{} string]}
& without it the result is \texttt{number[\allowbreak{}]}, then
\texttt{lo.\allowbreak{}toFixed(\allowbreak{})}:
\texttt{TS18048:\allowbreak{} 'lo1' is possibly 'undefined'.\allowbreak{}} \\
2
\texttt{const [\allowbreak{},\allowbreak{} x,\allowbreak{} .\allowbreak{}.\allowbreak{}.\allowbreak{}rest] =\allowbreak{} list}
& pattern, rest is an array & nothing, but \texttt{x} is
\texttt{T |\allowbreak{} undefined} &
\texttt{firstStop.\allowbreak{}toUpperCase(\allowbreak{})}:
\texttt{TS18048:\allowbreak{} 'firstStop' is possibly 'undefined'.\allowbreak{}} \\
3
\texttt{[\allowbreak{}a,\allowbreak{} b] =\allowbreak{} [\allowbreak{}b,\allowbreak{} a]}
& the swap & both variables must have the same type &
\texttt{let a:\allowbreak{} string,\allowbreak{} b:\allowbreak{} number}:
\texttt{TS2322:\allowbreak{} Type 'number' is not assignable to type 'string'.\allowbreak{}} \\
4 \texttt{const [\allowbreak{}a =\allowbreak{} 1] =\allowbreak{} list} &
defaults & the default must fit the item type &
\texttt{const [\allowbreak{}x =\allowbreak{} 5]:\allowbreak{} string[\allowbreak{}] =\allowbreak{} [\allowbreak{}]}:
\texttt{TS2322:\allowbreak{} Type 'number' is not assignable to type 'string'.\allowbreak{}} \\
5 \texttt{const \{ a \} =\allowbreak{} t} & the pattern & nothing, the
types come from \texttt{t} & a name not in the type:
\texttt{TS2339:\allowbreak{} Property 'platform' does not exist on type .\allowbreak{}.\allowbreak{}.\allowbreak{}} \\
6 \texttt{const \{ a:\allowbreak{} b \} =\allowbreak{} t} & rename & the
type goes after the pattern, not after the colon &
\texttt{\{ line:\allowbreak{} string,\allowbreak{} departure:\allowbreak{} string \}}:
\texttt{TS2451:\allowbreak{} Cannot redeclare block-scoped variable 'string'.\allowbreak{}} \\
7 \texttt{const \{ a =\allowbreak{} 1 \} =\allowbreak{} t} & defaults &
\texttt{?} on optional fields & no default, then
\texttt{const e:\allowbreak{} boolean =\allowbreak{} express}:
\texttt{TS2322:\allowbreak{} Type 'boolean |\allowbreak{} undefined' is not assignable to type 'boolean'.\allowbreak{}} \\
8
\texttt{const \{ a,\allowbreak{} .\allowbreak{}.\allowbreak{}.\allowbreak{}rest \} =\allowbreak{} t}
& rest & nothing, the rest type is computed
(\texttt{Omit<T,\allowbreak{} 'a'>}) &
\texttt{info.\allowbreak{}operator}:
\texttt{TS2339:\allowbreak{} Property 'operator' does not exist on type '\{ line:\allowbreak{} string; \}'.\allowbreak{}} \\
9 \texttt{const \{ p:\allowbreak{} \{ x \} \} =\allowbreak{} t} &
nesting & a default for an optional inner object &
\texttt{pier?:\allowbreak{} Pier} without default:
\texttt{TS2339:\allowbreak{} Property 'name' does not exist on type 'Pier |\allowbreak{} undefined'.\allowbreak{}} \\
10
\texttt{const \{ [\allowbreak{}k]:\allowbreak{} v \} =\allowbreak{} t} &
computed key & a \texttt{Record<string,\allowbreak{} T>} or
\texttt{keyof} type &
\texttt{v.\allowbreak{}toUpperCase(\allowbreak{})}:
\texttt{TS18048:\allowbreak{} 'v10' is possibly 'undefined'.\allowbreak{}} \\
11 \texttt{function f(\allowbreak{}\{ a \} =\allowbreak{} \{\})} &
pattern and defaults & a type after the pattern:
\texttt{(\allowbreak{}\{ a \}:\allowbreak{} Options =\allowbreak{} \{\})}
& no type:
\texttt{TS7031:\allowbreak{} Binding element 'line' implicitly has an 'any' type.\allowbreak{}} \\
12
\texttt{for (\allowbreak{}const [\allowbreak{}k,\allowbreak{} v] of map)}
& the loop &
\texttt{new Map<string,\allowbreak{} number>(\allowbreak{})} &
\texttt{count.\allowbreak{}toUpperCase(\allowbreak{})}:
\texttt{TS2339:\allowbreak{} Property 'toUpperCase' does not exist on type 'number'.\allowbreak{}} \\
\end{longtable}

\hypertarget{bonus-1-latex}{%
\section{Bonus 1: LaTeX}\label{bonus-1-latex}}

Bonus 1: LaTeX

The whole guide (every section of this README) is also a LaTeX document,
in the folder \texttt{latex/\allowbreak{}}:

\begin{itemize}
\tightlist
\item
  \texttt{latex/\allowbreak{}destructuring-guide.\allowbreak{}tex}: the
  LaTeX source.
\item
  \texttt{latex/\allowbreak{}destructuring-guide.\allowbreak{}pdf}: the
  PDF, with a title page, a linked table of contents, numbered sections,
  headers with my name, code typeset with the minted package (Pygments
  colours, line breaking for long lines) and the tables as real
  \texttt{longtable} tables.
\item
  \texttt{latex/\allowbreak{}build.\allowbreak{}sh},
  \texttt{latex/\allowbreak{}header.\allowbreak{}tex},
  \texttt{latex/\allowbreak{}minted.\allowbreak{}lua},
  \texttt{latex/\allowbreak{}metadata.\allowbreak{}yaml}: how it is
  made.
\end{itemize}

Commands I used (in a Linux shell with pandoc 3.1.3, TeX Live XeLaTeX
and Pygments 2.21):

\begin{minted}{text}
cd latex
sh build.sh
\end{minted}

\texttt{build.\allowbreak{}sh} does three things: it converts
\texttt{.\allowbreak{}.\allowbreak{}/\allowbreak{}README.\allowbreak{}md}
to LaTeX with pandoc, using \texttt{minted.\allowbreak{}lua} (a small
pandoc filter that turns every code block into a \texttt{minted}
environment) and \texttt{header.\allowbreak{}tex} (fonts, colours, page
layout); it gives the two wide tables fixed column widths so they wrap
inside the page; and it compiles with
\texttt{latexmk -xelatex -shell-escape destructuring-guide.\allowbreak{}tex}
(\texttt{-shell-escape} is needed by minted to call Pygments). The log
has no errors and no overfull boxes.

Why minted and not listings: I tried \texttt{listings} first, but under
XeLaTeX it cannot handle characters above code 255, and Turkish has
several (\texttt{ı}, \texttt{ş}, \texttt{ğ}, \texttt{İ}). The output was
\texttt{'ıKadköy-Karaköy-Eminönü'} instead of
\texttt{'Kadıköy-Karaköy-Eminönü'}: listings moved the \texttt{ı} to the
start of the word. Even the usual fix
(\texttt{literate=\allowbreak{}\{ı\}\{\{\textbackslash{}i\}\}1 .\allowbreak{}.\allowbreak{}.\allowbreak{}})
did not change it. minted gives the text to Pygments and back to XeLaTeX
as Unicode, so every letter is right. The one character I had to replace
is the ship emoji in D1, because the DejaVu fonts have no colour emoji;
in the PDF it is written \texttt{[\allowbreak{}ship]}.

\hypertarget{bonus-2-depth}{%
\section{Bonus 2: Depth}\label{bonus-2-depth}}

Bonus 2: Depth

Code: \texttt{js/\allowbreak{}depth.\allowbreak{}mjs} (D1--D4) and
\texttt{ts/\allowbreak{}depth-types.\allowbreak{}tsx} (D5--D8). Run with
\texttt{npm run bonus}. The outputs below are from that run.

\hypertarget{d1-what-an-array-pattern-really-calls}{%
\subsection{D1: what an array pattern really
calls}\label{d1-what-an-array-pattern-really-calls}}

\begin{minted}{javascript}
const [k, a, d] = 'Kadıköy'
const [ship, firstLetter] = '[ship]Kadıköy'
const [p1, p2, p3] = new Set(['Kadıköy', 'Kadıköy', 'Eminönü'])
const [[firstPier, firstCount]] = new Map([['Kadıköy', 5], ['Beşiktaş', 2]])
\end{minted}

\begin{minted}{text}
D1 string: K a d
D1 emoji by pattern: [ship] K
D1 emoji by index: "\ud83d" "\udea2" (two halves of one emoji)
D1 Set: Kadıköy Eminönü undefined
D1 Map first entry: Kadıköy 5
\end{minted}

An array pattern does not read \texttt{value[\allowbreak{}0]},
\texttt{value[\allowbreak{}1]}. It calls
\texttt{value[\allowbreak{}Symbol.\allowbreak{}iterator](\allowbreak{})}
and then \texttt{next(\allowbreak{})} once for each place in the
pattern. That is why it works on strings, \texttt{Set}, \texttt{Map} and
generators. It also explains the emoji: a string\textquotesingle s
iterator walks by code points, so the ship comes out whole, while
\texttt{shipLine[\allowbreak{}0]} reads one UTF-16 unit, which is only
half of the emoji. A \texttt{Set} iterator never gives the duplicate, so
the third place is \texttt{undefined}.

To see the calls, I used a generator that logs every pier it gives and
has a \texttt{finally} block:

\begin{minted}{javascript}
function* halicPiers() {
  try {
    for (const pier of ['Beşiktaş', 'Kabataş', 'Karaköy', 'Kasımpaşa']) {
      console.log('D1   generator gives:', pier)
      yield pier
    }
    console.log('D1   generator loop ended normally')
  } finally {
    console.log('D1   generator finally block')
  }
}
const [g1, g2] = halicPiers()
const [, ...afterFirst] = halicPiers()
\end{minted}

\begin{minted}{text}
D1   generator gives: Beşiktaş
D1   generator gives: Kabataş
D1   generator finally block
D1 generator result: Beşiktaş Kabataş
D1   generator gives: Beşiktaş
D1   generator gives: Kabataş
D1   generator gives: Karaköy
D1   generator gives: Kasımpaşa
D1   generator loop ended normally
D1   generator finally block
D1 generator rest: [ 'Kabataş', 'Karaköy', 'Kasımpaşa' ]
\end{minted}

With \texttt{[\allowbreak{}g1,\allowbreak{} g2]} the generator gave only
two piers and then its \texttt{finally} ran without "loop ended
normally". So the pattern did not run the generator to the end; it took
two values and then closed the iterator early by calling its
\texttt{return(\allowbreak{})} method (the specification calls this
\emph{IteratorClose}). With a rest element the pattern must ask for
everything, so the loop really ends. A first version of this experiment
printed "return() was called" in the \texttt{finally} block, which was
wrong for the rest case: \texttt{finally} also runs when the loop ends
normally. The extra "loop ended normally" line is what tells the two
cases apart.

A custom object with only \texttt{Symbol.\allowbreak{}iterator} works
too, and \texttt{[\allowbreak{}t1]} called \texttt{next(\allowbreak{})}
exactly once (\texttt{D1   next(\allowbreak{}) call 1}): the pattern
does not ask for one more item to check whether the iterator is
finished. A plain object has no \texttt{Symbol.\allowbreak{}iterator},
so
\texttt{const [\allowbreak{}x] =\allowbreak{} \{ line:\allowbreak{} 'Kadıköy-Beşiktaş' \}}
is
\texttt{TypeError - \{(\allowbreak{}intermediate value)\} is not iterable}.

\hypertarget{d2-when-defaults-are-evaluated}{%
\subsection{D2: when defaults are
evaluated}\label{d2-when-defaults-are-evaluated}}

\begin{minted}{javascript}
function fallbackPier(reason) {
  console.log(`D2   fallbackPier() called (${reason})`)
  return 'Eminönü'
}
const [given = fallbackPier('first pattern')] = ['Kadıköy']
const [missing = fallbackPier('second pattern')] = []
const [from = 'Kadıköy', to = from] = []
new Function('const [x = y, y = 1] = []; return x')()
\end{minted}

\begin{minted}{text}
D2 value given: Kadıköy
D2   fallbackPier() called (second pattern)
D2 value missing: Eminönü
D2 earlier variable: Kadıköy Kadıköy
D2 later variable: ReferenceError - Cannot access 'y' before initialization
\end{minted}

A default is an expression that runs only when it is needed: the
function with a side effect was not called for the first pattern at all.
Defaults run left to right, one variable after the other, so a default
can use a variable that is \emph{earlier} in the same pattern
(\texttt{to =\allowbreak{} from}), but a \emph{later} one is still in
the temporal dead zone, the same as using a \texttt{let} before its
line.

The order of reading also follows the pattern, not the object. With
getters that log:

\begin{minted}{javascript}
const watched = {
  get line() { console.log('D2   read line'); return 'Kadıköy-Beşiktaş' },
  get departure() { console.log('D2   read departure'); return '18:10' },
}
const { departure: dep, line: lin } = watched
\end{minted}

\begin{minted}{text}
D2   read departure
D2   read line
\end{minted}

\texttt{departure} is read first because it is first in my pattern,
although \texttt{line} is first in the object.

\hypertarget{d3-destructuring-in-other-places}{%
\subsection{D3: destructuring in other
places}\label{d3-destructuring-in-other-places}}

\begin{minted}{javascript}
try { null.line } catch ({ name, message }) { console.log('D3 catch:', name, '|', message) }

for (const [firstChar, secondChar] in counts) { ... }          // counts = { Kadıköy: 5, Beşiktaş: 2 }
for (const [pier, count] of Object.entries(counts)) { ... }

;[board.from, board.to] = ['Kabataş', 'Kadıköy']
;[route[0], route[2]] = [route[2], route[0]]
;({ line: board.line = 'bilinmiyor', operator: board.operator } = { operator: 'Şehir Hatları' })

const { renderToStaticMarkup: render } = await import('react-dom/server')
\end{minted}

\begin{minted}{text}
D3 catch: TypeError | Cannot read properties of null (reading 'line')
D3 for...in pattern: K a
D3 for...in pattern: B e
D3 for...of entries: Kadıköy 5
D3 for...of entries: Beşiktaş 2
D3 property targets: { from: 'Kabataş', to: 'Kadıköy' }
D3 swap array items: [ 'Kadıköy', 'Karaköy', 'Eminönü' ]
D3 object pattern into properties, with default: { from: 'Kabataş', to: 'Kadıköy', line: 'bilinmiyor', operator: 'Şehir Hatları' }
D3 dynamic import + rename: function
\end{minted}

\begin{itemize}
\tightlist
\item
  \texttt{catch} takes a pattern like a parameter does (form 11), so I
  can take only \texttt{name} and \texttt{message}.
\item
  \texttt{for.\allowbreak{}.\allowbreak{}.\allowbreak{}in} gives the
  keys as strings. A pattern there is legal but means something
  surprising: \texttt{[\allowbreak{}firstChar,\allowbreak{} secondChar]}
  takes the first two letters of each key (a string is iterable, D1).
  For pairs you need
  \texttt{for.\allowbreak{}.\allowbreak{}.\allowbreak{}of} with
  \texttt{Object.\allowbreak{}entries} (form 12).
\item
  In an assignment (not a declaration) the targets can be any place you
  can assign to: object properties, array items, even with defaults.
  Swapping two items of an array is form 3 with
  \texttt{route[\allowbreak{}0]} and \texttt{route[\allowbreak{}2]} as
  targets.
\item
  \texttt{import \{ a as b \} from '.\allowbreak{}.\allowbreak{}.\allowbreak{}'}
  looks like destructuring but is a different syntax: it renames with
  \texttt{as}, not \texttt{:\allowbreak{}}, it cannot have defaults and
  it cannot be nested. A \emph{dynamic} \texttt{import(\allowbreak{})}
  returns a normal object (a promise of it), so there a real object
  pattern with a rename works.
\end{itemize}

\hypertarget{d4-when-destructuring-makes-code-worse}{%
\subsection{D4: when destructuring makes code
worse}\label{d4-when-destructuring-makes-code-worse}}

\begin{minted}{javascript}
const { sailing: { pier: { name: deepName }, times: { first: deepFirst, last: deepLast } } } = trip
\end{minted}

\begin{minted}{text}
D4 deep pattern: Kadıköy 07:30 23:45
\end{minted}

It works, but it is hard to read (three levels, three renames on one
line), and if any level is missing the whole line throws (form 9).
Refactored, one level per line and with \texttt{?.\allowbreak{}} where
data can be missing:

\begin{minted}{javascript}
const { pier, times } = trip.sailing
const pierName = pier?.name ?? 'bilinmiyor'
const { first, last } = times
\end{minted}

\begin{minted}{text}
D4 refactored: Kadıköy 07:30 23:45
\end{minted}

My rule after this: at most one nested level in a pattern, and no rename
inside a nested part. Long lists of renames are a sign that a plain
\texttt{const x =\allowbreak{} obj.\allowbreak{}y} is clearer.

\hypertarget{d5-infer-is-a-pattern-for-types}{%
\subsection{\texorpdfstring{D5: \texttt{infer} is a pattern for
types}{D5:  is a pattern for types}}\label{d5-infer-is-a-pattern-for-types}}

\begin{minted}{typescript}
type First<T extends readonly unknown[]> = T extends readonly [infer F, ...unknown[]] ? F : undefined
type Rest<T extends readonly unknown[]> = T extends readonly [unknown, ...infer R] ? R : []
type Last<T extends readonly unknown[]> = T extends readonly [...unknown[], infer L] ? L : undefined

const kadikoyLine = ['Kadıköy', 'Karaköy', 'Eminönü'] as const
const start: First<typeof kadikoyLine> = 'Kadıköy'
const others: Rest<typeof kadikoyLine> = ['Karaköy', 'Eminönü']
const end: Last<typeof kadikoyLine> = 'Eminönü'

const [startValue, ...othersValue] = kadikoyLine
const sameStart: First<typeof kadikoyLine> = startValue
const sameOthers: Rest<typeof kadikoyLine> = othersValue
\end{minted}

\begin{minted}{text}
D5 First/Rest/Last: Kadıköy [ 'Karaköy', 'Eminönü' ] Eminönü
D5 same as the value pattern: Kadıköy [ 'Karaköy', 'Eminönü' ]
D5 First<string[]>: undefined
\end{minted}

\texttt{T extends [\allowbreak{}infer F,\allowbreak{} .\allowbreak{}.\allowbreak{}.\allowbreak{}infer R]}
is the type-level twin of
\texttt{const [\allowbreak{}f,\allowbreak{} .\allowbreak{}.\allowbreak{}.\allowbreak{}r] =\allowbreak{} value}:
TypeScript tries to match the tuple type against the pattern and, if it
fits, gives names to the parts. \texttt{Last} shows something the value
level cannot do: a rest element before the last one
(\texttt{[\allowbreak{}.\allowbreak{}.\allowbreak{}.\allowbreak{}unknown[\allowbreak{}],\allowbreak{} infer L]})
is allowed in types, but
\texttt{const [\allowbreak{}.\allowbreak{}.\allowbreak{}.\allowbreak{}a,\allowbreak{} last] =\allowbreak{} list}
is a syntax error. The checks are real:
\texttt{const wrong:\allowbreak{} First<typeof l> =\allowbreak{} 'Karaköy'}
gives
\texttt{error TS2322:\allowbreak{} Type '"Karaköy"' is not assignable to type '"Kadıköy"'.\allowbreak{}}
\texttt{First<string[\allowbreak{}]>} is \texttt{undefined}, because a
\texttt{string[\allowbreak{}]} does not match a tuple with at least one
item; this is the type-level version of T3.

\hypertarget{d6d8-objects-parameters-and-generic-components}{%
\subsection{D6--D8: objects, parameters and generic
components}\label{d6d8-objects-parameters-and-generic-components}}

\begin{minted}{typescript}
type PierNameOf<T> = T extends { pier: { name: infer N } } ? N : never   // D6: like { pier: { name } }
type DescribeOptions = Parameters<typeof describeSailing>[0]             // D7: the form 11 options type
type PierBadgeProps = ComponentProps<typeof PierBadge>                   // D7: props of my component

type ListProps<T> = { items: readonly T[]; children: (item: T, index: number) => ReactNode }
function List<T>({ items, children }: ListProps<T>) { ... }              // D8
<List items={trips}>{({ from, to, minutes }) => `${from} → ${to}, ${minutes} dk`}</List>
\end{minted}

\begin{minted}{text}
D6 nested infer: Beşiktaş
D7 Parameters: Kadıköy-Beşiktaş (ekspres)
D7 ComponentProps<typeof>: <span>Üsküdar (Anadolu)</span>
D8 generic component: <ul><li>Kabataş → Kadıköy, 20 dk</li><li>Beşiktaş → Üsküdar, 10 dk</li></ul>
\end{minted}

\texttt{infer} works inside object types the same way (D6 is form 9 for
types). \texttt{Parameters<F>} takes a function type apart into a tuple,
and \texttt{[\allowbreak{}0]} is the options object of form 11, so a
caller can reuse the type without exporting it. In D8 the component is
generic: \texttt{T} is inferred from \texttt{items}, so the render prop
(P3) gets \texttt{\{ from,\allowbreak{} to,\allowbreak{} minutes \}}
with the type \texttt{Trip} without my writing it. When I gave items
without \texttt{to}, the pattern \texttt{(\allowbreak{}\{ to \})} was
rejected:
\texttt{error TS2339:\allowbreak{} Property 'to' does not exist on type '\{ from:\allowbreak{} string; \}'.\allowbreak{}}

\hypertarget{bonus-3-real-code}{%
\section{Bonus 3: Real code}\label{bonus-3-real-code}}

Bonus 3: Real code

Two sources, both with exact file and line:

\begin{itemize}
\tightlist
\item
  \textbf{A library I have installed:} \texttt{esbuild} 0.28.2 in
  \texttt{node\_modules/\allowbreak{}esbuild/\allowbreak{}lib/\allowbreak{}main.\allowbreak{}js}
  (it is installed by \texttt{tsx}, which runs every example in this
  project).
\item
  \textbf{An open-source project:}
  \href{https://github.com/pacocoursey/cmdk}{cmdk}, a React command-menu
  component, file
  \texttt{cmdk/\allowbreak{}src/\allowbreak{}index.\allowbreak{}tsx} at
  commit
  \href{https://github.com/pacocoursey/cmdk/blob/dd2250ed608443e8f32bafc5fa2d1d07a3746aa3/cmdk/src/index.tsx}{\texttt{dd2250e}}
  (2025-10-28). Line numbers are for that commit.
\end{itemize}

\hypertarget{case-1-esbuild-mainjs193-form-5}{%
\subsection{\texorpdfstring{Case 1: esbuild,
\texttt{main.\allowbreak{}js:\allowbreak{}193} (form
5)}{Case 1: esbuild,  (form 5)}}\label{case-1-esbuild-mainjs193-form-5}}

\begin{minted}{javascript}
decodeUTF8 = (bytes) => {
  let { buffer, byteOffset, byteLength } = bytes;
  return Buffer.from(buffer, byteOffset, byteLength).toString();
};
\end{minted}

\textbf{Without destructuring:}
\texttt{return Buffer.\allowbreak{}from(\allowbreak{}bytes.\allowbreak{}buffer,\allowbreak{} bytes.\allowbreak{}byteOffset,\allowbreak{} bytes.\allowbreak{}byteLength).\allowbreak{}toString(\allowbreak{});}

Form 5: three properties by name. The variable names are exactly the
names of the parameters of \texttt{Buffer.\allowbreak{}from}, so the
pattern reads like documentation of what the call needs. Here I think it
is a good use.

\hypertarget{case-2-esbuild-mainjs17451747-form-5-and-pitfall-1-avoided}{%
\subsection{\texorpdfstring{Case 2: esbuild,
\texttt{main.\allowbreak{}js:\allowbreak{}1745–1747} (form 5, and
pitfall 1
avoided)}{Case 2: esbuild,  (form 5, and pitfall 1 avoided)}}\label{case-2-esbuild-mainjs17451747-form-5-and-pitfall-1-avoided}}

\begin{minted}{javascript}
if (!e.location) return `
error: ${e.text}`;
let { file, line, column } = e.location;
\end{minted}

Form 5 again, but notice the line before it:
\texttt{e.\allowbreak{}location} can be missing, and
\texttt{let \{ file \} =\allowbreak{} undefined} would be the
\texttt{TypeError} of pitfall 1. The authors check first and return
early. \textbf{Without destructuring:}
\texttt{const file =\allowbreak{} e.\allowbreak{}location.\allowbreak{}file}
(and the same check is still needed).

\hypertarget{case-3-esbuild-mainjs2093-form-1-on-a-string-array}{%
\subsection{\texorpdfstring{Case 3: esbuild,
\texttt{main.\allowbreak{}js:\allowbreak{}2093} (form 1 on a string
array)}{Case 3: esbuild,  (form 1 on a string array)}}\label{case-3-esbuild-mainjs2093-form-1-on-a-string-array}}

\begin{minted}{javascript}
let [major, minor] = process.versions.node.split(".");
if (+major < 12 || +major === 12 && +minor < 17 || +major === 13 && +minor < 13) {
\end{minted}

Form 1: values by position from
\texttt{"22.\allowbreak{}23.\allowbreak{}4".\allowbreak{}split(\allowbreak{}".\allowbreak{}")}.
\textbf{Without destructuring:}
\texttt{const parts =\allowbreak{} process.\allowbreak{}versions.\allowbreak{}node.\allowbreak{}split(\allowbreak{}".\allowbreak{}"); const major =\allowbreak{} parts[\allowbreak{}0]; const minor =\allowbreak{} parts[\allowbreak{}1];}.
The patch number is simply not named (the pattern is shorter than the
array, and that is fine). A detail: \texttt{major} and \texttt{minor}
are strings, which is why every comparison has a \texttt{+} in front to
make them numbers. In TypeScript with \texttt{noUncheckedIndexedAccess}
they would be \texttt{string |\allowbreak{} undefined} (T3).

\hypertarget{case-4-esbuild-mainjs2271-form-1-a-function-that-returns-a-pair}{%
\subsection{\texorpdfstring{Case 4: esbuild,
\texttt{main.\allowbreak{}js:\allowbreak{}2271} (form 1, a function that
returns a
pair)}{Case 4: esbuild,  (form 1, a function that returns a pair)}}\label{case-4-esbuild-mainjs2271-form-1-a-function-that-returns-a-pair}}

\begin{minted}{javascript}
let [command, args] = esbuildCommandAndArgs();
let child = child_process.spawn(command, args.concat(`--service=${"0.28.2"}`, "--ping"), { ... });
\end{minted}

Exactly my form 1 (\texttt{nextFerry(\allowbreak{})}): a function
returns two values that belong together. \textbf{Without destructuring:}
\texttt{const result =\allowbreak{} esbuildCommandAndArgs(\allowbreak{}); const command =\allowbreak{} result[\allowbreak{}0]; const args =\allowbreak{} result[\allowbreak{}1];}.
The same line appears again at
\texttt{main.\allowbreak{}js:\allowbreak{}2364}.

\hypertarget{case-5-esbuild-mainjs2394-form-6-rename}{%
\subsection{\texorpdfstring{Case 5: esbuild,
\texttt{main.\allowbreak{}js:\allowbreak{}2394} (form 6,
rename)}{Case 5: esbuild,  (form 6, rename)}}\label{case-5-esbuild-mainjs2394-form-6-rename}}

\begin{minted}{javascript}
let { port1: mainPort, port2: workerPort } = new worker_threads2.MessageChannel();
\end{minted}

Form 6: \texttt{MessageChannel} has two ports with meaningless names
(\texttt{port1}, \texttt{port2}); the rename gives each one the role it
has in this code. \textbf{Without destructuring:}
\texttt{const channel =\allowbreak{} new MessageChannel(\allowbreak{}); const mainPort =\allowbreak{} channel.\allowbreak{}port1; const workerPort =\allowbreak{} channel.\allowbreak{}port2;}.
This is the best "use it when" for form 6 that I found.

\hypertarget{case-6-cmdk-command-indextsx191202-forms-5-7-and-8-on-react-props}{%
\subsection{\texorpdfstring{Case 6: cmdk \texttt{Command},
\texttt{index.\allowbreak{}tsx:\allowbreak{}191–202} (forms 5, 7 and 8
on React
props)}{Case 6: cmdk ,  (forms 5, 7 and 8 on React props)}}\label{case-6-cmdk-command-indextsx191202-forms-5-7-and-8-on-react-props}}

\begin{minted}{tsx}
const {
  label,
  children,
  value,
  onValueChange,
  filter,
  shouldFilter,
  loop,
  disablePointerSelection = false,
  vimBindings = true,
  ...etc
} = props
\end{minted}

and later, line 577:
\texttt{<Primitive.\allowbreak{}div .\allowbreak{}.\allowbreak{}.\allowbreak{} \{.\allowbreak{}.\allowbreak{}.\allowbreak{}etc\}>}.

This is P4 in a real library: the props the component understands are
taken out by name (form 5), two of them get defaults (form 7, like P6),
and everything else goes into \texttt{etc} (form 8) and is spread onto
the real \texttt{<div>}, so a user can pass \texttt{className},
\texttt{style} or \texttt{data-*} attributes. \textbf{Props and children
here:} \texttt{label}, \texttt{value}, \texttt{filter}, \texttt{loop}
... are props (attributes the user writes in JSX); \texttt{children} is
the menu content written between \texttt{<Command>} and
\texttt{</\allowbreak{}Command>}. \textbf{Without destructuring:} eleven
lines of \texttt{const label =\allowbreak{} props.\allowbreak{}label},
two with
\texttt{=\allowbreak{}=\allowbreak{}=\allowbreak{} undefined ? .\allowbreak{}.\allowbreak{}.\allowbreak{} :\allowbreak{} .\allowbreak{}.\allowbreak{}.\allowbreak{}},
and a manual copy of \texttt{props} with nine \texttt{delete}s for
\texttt{etc}.

\hypertarget{case-7-cmdk-group-indextsx730-props-vs-children-p5}{%
\subsection{\texorpdfstring{Case 7: cmdk \texttt{Group},
\texttt{index.\allowbreak{}tsx:\allowbreak{}730} (props vs children,
P5)}{Case 7: cmdk ,  (props vs children, P5)}}\label{case-7-cmdk-group-indextsx730-props-vs-children-p5}}

\begin{minted}{tsx}
const { heading, children, forceMount, ...etc } = props
...
{heading && (
  <div ref={headingRef} cmdk-group-heading="" aria-hidden id={headingId}>
    {heading}
  </div>
)}
\end{minted}

This is exactly my P5 decision. A group has a heading and items. The
heading is a \textbf{prop}, because the component must know which part
is the heading: it puts it in its own \texttt{<div>}, gives it an
\texttt{id} and \texttt{aria-hidden}, and does not render it at all when
it is missing
(\texttt{heading \&\& .\allowbreak{}.\allowbreak{}.\allowbreak{}}, the
\texttt{false} case of P3). The items are children, because the group
does not need to understand them, it only wraps them. Used like:
\texttt{<Command.\allowbreak{}Group heading=\allowbreak{}"Piers"><Command.\allowbreak{}Item>Kadıköy</\allowbreak{}Command.\allowbreak{}Item></\allowbreak{}Command.\allowbreak{}Group>}.

\textbf{Without destructuring:}
\texttt{const heading =\allowbreak{} props.\allowbreak{}heading},
\texttt{const children =\allowbreak{} props.\allowbreak{}children},
\texttt{const forceMount =\allowbreak{} props.\allowbreak{}forceMount},
and \texttt{etc} as a copy of \texttt{props} with those three deleted.
The code even mixes both styles: it destructures \texttt{heading} but
later still writes \texttt{props.\allowbreak{}heading} (line 744,
\texttt{useValue(\allowbreak{}id,\allowbreak{} ref,\allowbreak{} [\allowbreak{}props.\allowbreak{}value,\allowbreak{} props.\allowbreak{}heading,\allowbreak{} headingRef])}).

\hypertarget{case-8-cmdk-list-indextsx834-a-default-for-a-prop-p6}{%
\subsection{\texorpdfstring{Case 8: cmdk \texttt{List},
\texttt{index.\allowbreak{}tsx:\allowbreak{}834} (a default for a prop,
P6)}{Case 8: cmdk ,  (a default for a prop, P6)}}\label{case-8-cmdk-list-indextsx834-a-default-for-a-prop-p6}}

\begin{minted}{tsx}
const { children, label = 'Suggestions', ...etc } = props
\end{minted}

\texttt{label} is used for \texttt{aria-label} of the list. The default
only applies when the prop is \texttt{undefined}, so
\texttt{<Command.\allowbreak{}List label=\allowbreak{}\{null\}>} would
give no label at all (P6 and pitfall 2).

\textbf{Without destructuring:}
\texttt{const label =\allowbreak{} props.\allowbreak{}label =\allowbreak{}=\allowbreak{}=\allowbreak{} undefined ? 'Suggestions' :\allowbreak{} props.\allowbreak{}label}
(plus \texttt{children} and the copy for \texttt{etc}). The pattern
makes the default visible right where the props are read.

\hypertarget{case-9-where-destructuring-makes-the-code-worse-cmdk-item-indextsx704}{%
\subsection{\texorpdfstring{Case 9 (where destructuring makes the code
worse): cmdk \texttt{Item},
\texttt{index.\allowbreak{}tsx:\allowbreak{}704}}{Case 9 (where destructuring makes the code worse): cmdk , }}\label{case-9-where-destructuring-makes-the-code-worse-cmdk-item-indextsx704}}

\begin{minted}{tsx}
const { disabled, value: _, onSelect: __, forceMount: ___, keywords: ____, ...etc } = props
\end{minted}

The goal is to put "every prop except these four" into \texttt{etc}, so
that \texttt{value}, \texttt{onSelect}, \texttt{forceMount} and
\texttt{keywords} are not spread onto the \texttt{<div>} as HTML
attributes. Because the only way to leave a property out of an object
rest is to name it in the pattern, the authors rename four props to the
throw-away names \texttt{\_}, \texttt{\_\_}, \texttt{\_\_\_},
\texttt{\_\_\_\_} (form 6 used only to avoid a name clash). The names
say nothing, a fifth excluded prop would need \texttt{\_\_\_\_\_}, and a
reader has to know the trick to understand why the line is there.
\textbf{Without destructuring for the exclusion,} a small helper is
clearer:

\begin{minted}{typescript}
function omit<T extends object, K extends keyof T>(obj: T, ...keys: K[]): Omit<T, K> {
  const copy = { ...obj }
  for (const key of keys) delete copy[key]
  return copy
}
const etc = omit(props, 'disabled', 'value', 'onSelect', 'forceMount', 'keywords')
\end{minted}

I checked that this helper compiles in strict mode; with
\texttt{\{ disabled:\allowbreak{} true,\allowbreak{} value:\allowbreak{} 'x',\allowbreak{} className:\allowbreak{} 'c' \}}
it returns \texttt{\{ className:\allowbreak{} 'c' \}}. It is longer, but
it says what it does ("omit these") instead of hiding it in renames.

(The handout also suggests the Task Manager starter from Assignment 3;
it is not in this project folder, so my cases come from the installed
library and the open-source project.)

\hypertarget{my-own-log}{%
\section{My own log}\label{my-own-log}}

Things that really went wrong or surprised me while I worked on this
guide, with the exact messages.

\begin{enumerate}
\def\labelenumi{\arabic{enumi}.}
\item
  \textbf{I ran \texttt{npm install} in the wrong folder.} I opened the
  terminal in \texttt{webproje} and not in
  \texttt{webproje\textbackslash{}assignment4}, and got:
  \texttt{npm error enoent Could not read package.\allowbreak{}json:\allowbreak{} Error:\allowbreak{} ENOENT:\allowbreak{} no such file or directory,\allowbreak{} open 'C:\allowbreak{}\textbackslash{}Users\textbackslash{}Şevval\textbackslash{}webproje\textbackslash{}package.\allowbreak{}json'}.
  npm looks for \texttt{package.\allowbreak{}json} only in the folder
  the terminal is in. After \texttt{cd assignment4} it worked. Since
  then I open the \texttt{assignment4} folder itself in VS Code, so the
  terminal starts in the right place.
\item
  \textbf{My wrong guess about form 2 with a one-pier line.} Before
  running
  \texttt{const [\allowbreak{},\allowbreak{} onlyStop,\allowbreak{} .\allowbreak{}.\allowbreak{}.\allowbreak{}afterThat] =\allowbreak{} [\allowbreak{}'Kabataş']}
  I wrote my guess in the file:
  \texttt{onlyStop hata verir,\allowbreak{} afterThat boş dizi olur}
  ("onlyStop throws an error, afterThat is an empty array"). The real
  output was
  \texttt{form 2 (\allowbreak{}one pier):\allowbreak{} undefined [\allowbreak{}]}.
  I was right about the rest, but wrong about \texttt{onlyStop}: reading
  an item that is not there is not an error in JavaScript, it is just
  \texttt{undefined}, the same as \texttt{shortLine[\allowbreak{}1]}.
\item
  \textbf{My wrong guess about form 9 with a missing \texttt{pier}.}
  This time I guessed the opposite: after form 2 I thought a missing
  part always gives \texttt{undefined}, so \texttt{nightPier} would be
  \texttt{undefined} and the program would continue. The real output was
  \texttt{Cannot read properties of undefined (\allowbreak{}reading 'name')}.
  Writing the pattern as two plain lines showed me why:
  \texttt{const pier =\allowbreak{} nightSailing.\allowbreak{}pier} is
  \texttt{undefined} (fine), but \texttt{pier.\allowbreak{}name} tries
  to go \emph{into} \texttt{undefined}, and that throws. Reading a
  missing value is OK, going inside it is not.
\item
  \textbf{TypeScript did not let me call a method on a destructured
  item.} I changed \texttt{firstStop ?? 'no stop'} to
  \texttt{firstStop.\allowbreak{}toUpperCase(\allowbreak{})} in
  \texttt{ts/\allowbreak{}arrays.\allowbreak{}ts} and
  \texttt{npm run check} said:
  \texttt{ts/\allowbreak{}arrays.\allowbreak{}ts:\allowbreak{}16:\allowbreak{}24 - error TS18048:\allowbreak{} 'firstStop' is possibly 'undefined'.\allowbreak{}}
  It is the same \texttt{undefined} as in entry 2, but TypeScript finds
  it before the code runs, because \texttt{noUncheckedIndexedAccess} is
  on and a \texttt{string[\allowbreak{}]} can be shorter than the
  pattern. I put the \texttt{??} back.
\item
  \textbf{I pasted terminal output back into the terminal (twice).} I
  copied the output to show it, and then pasted it into the PowerShell
  terminal instead of my command. PowerShell tried to run every line:
  \texttt{The term 'form' is not recognized as the name of a cmdlet .\allowbreak{}.\allowbreak{}.\allowbreak{}},
  and for the lines that started with my prompt:
  \texttt{Get-Process :\allowbreak{} A positional parameter cannot be found that accepts argument 'npm'.\allowbreak{}}
  The second message confused me until I learned that in PowerShell
  \texttt{PS} is a short name (alias) for \texttt{Get-Process}, so
  \texttt{PS C:\allowbreak{}\textbackslash{}.\allowbreak{}.\allowbreak{}.\allowbreak{}> npm run check}
  was read as \texttt{Get-Process} with strange arguments. Nothing was
  changed, no command really ran. Now I type the command or copy only
  that one line.
\item
  \textbf{\texttt{npm warn install-scripts 1 package has install scripts not yet covered by allowScripts:\allowbreak{} esbuild@0.\allowbreak{}28.\allowbreak{}2 (\allowbreak{}postinstall:\allowbreak{} node install.\allowbreak{}js)}}
  My npm does not run install scripts of packages until I allow them. I
  thought \texttt{tsx} would not work without it, but
  \texttt{npm run all} worked, because esbuild\textquotesingle s Windows
  binary comes as a separate optional package
  (\texttt{@esbuild/\allowbreak{}win32-x64}) and does not need that
  script.
\item
  \textbf{\texttt{warning:\allowbreak{} in the working copy of 'js/\allowbreak{}arrays.\allowbreak{}mjs',\allowbreak{} LF will be replaced by CRLF the next time Git touches it}}
  I got this on almost every \texttt{git add}. It is not an error: the
  files have Linux line endings (LF) and Git on Windows will
  store/convert them to Windows line endings (CRLF). The code and the
  output do not change.
\item
  \textbf{\texttt{output.\allowbreak{}txt} and PowerShell.} The handout
  says \texttt{npm run all > output.\allowbreak{}txt}, but in Windows
  PowerShell \texttt{>} writes the file as UTF-16, which breaks how the
  file shows on GitHub. I used
  \texttt{cmd /\allowbreak{}c "npm run all > output.\allowbreak{}txt"}
  so that the file is UTF-8 and \texttt{Kadıköy} stays \texttt{Kadıköy}.
  Two React warnings still appeared in the terminal and not in the file,
  because they go to stderr and \texttt{>} only saves stdout.
\end{enumerate}

\end{document}
